% Version 2.20 of 2017/10/04

\documentclass[runningheads]{llncs}
\pagestyle{plain}
% !TEX root = main.tex
\usepackage{amsmath}
\let\proof\relax
\let\endproof\relax
\usepackage{amsthm}
\usepackage{amssymb}
\usepackage[linesnumbered,ruled,vlined]{algorithm2e}
\usepackage{booktabs} %toprule, bottomrule
\usepackage{bm}

%\usepackage{enumerate}
\usepackage{enumitem}

\usepackage{graphicx}
% Sort numerical citations within each group; retain range compression.
\usepackage[sort]{cite}
\usepackage[breaklinks=true]{hyperref}

\usepackage{makecell}
\usepackage{multirow}
\usepackage{mathtools}

\usepackage{stmaryrd}
% St Mary's Road symbols have no bold font; keep their regular shapes.
\SetSymbolFont{stmry}{bold}{U}{stmry}{m}{n}
\usepackage{threeparttable} % e.g. for tnote

\usepackage{xcolor}
\usepackage{xspace}
\usepackage{arydshln}

\usepackage{boxedminipage}
\usepackage[most]{tcolorbox}
\usepackage{mdframed}
%\usepackage{tikz}
%\usetikzlibrary{matrix,decorations.pathreplacing}

%\usepackage{showframe}
\usepackage{todonotes}

%\usepackage{amsfonts}
%\usepackage{float}
%\usepackage{lipsum}
%\usepackage{url}

%\usepackage[showframe]{geometry}

% used for tikz
\usetikzlibrary{positioning, shapes, arrows.meta, calc, backgrounds, fit, calc}
\usetikzlibrary{decorations.pathreplacing}
\renewcommand{\ttdefault}{lmtt} % use Latin Modern mono instead of cmtt

% !TEX root = main.tex
\newcommand{\methodname}{{\tt Mikros}\xspace}
\newcommand{\qafold}{{\tt QAFold}\xspace}
\newcommand{\zkSNARK}{{zkSNARK}\xspace}
\newcommand{\zkSNARKs}{{\tt zkSNARKs}\xspace}
\newcommand{\litebullet}{{\tt LiteBullet}\xspace}
\newcommand{\codeswitch}{{\tt CodeSwitch}\xspace}
% Define custom commands for Dirac notation
\newcommand{\bra}[1]{\langle #1|}
\newcommand{\ket}[1]{|#1\rangle}
\newcommand{\braket}[2]{\langle \, #1 \,|\,	#2 \,\rangle}
\newcommand{\dirac}[3]{\langle \, #1 \,|\,	#2 \,|\,	#3 \,\rangle}
\newcommand{\diracbra}[2]{\langle \, #1 \,|\,	#2 \,|}
\newcommand{\diracket}[2]{|\,	#2 \,|\,	#1 \,\rangle}
\newcommand{\diracbraket}[4]{\langle \, #1 \,|\,	#2 \,|\,	#3 \,|\, #4 \, \rangle}
\newcommand{\ExtensionDegree}[2]{[ \, #1 : #2 \,]}
\newcommand{\ext}[1]{\overline{#1}}
\ExplSyntaxOn
\NewDocumentCommand{\concat}{m}
 {
  \left[\, \clist_use:nn { #1 } { \mid } \,\right]
 }
\ExplSyntaxOff
\newcommand{\blockmat}[4]{
    \begin{pmatrix}
     #1 &\ \ \ & #2 \\
     #3 &\ \ \ & #4 \\
    \end{pmatrix}
}
\newcommand{\foldmat}[2]{
    \begin{pmatrix}
     \begin{pmatrix}
     #1 \\
     #1 \\
     \end{pmatrix}
     &
     \ \ 
     #2
    \begin{pmatrix}
     #1 \\
     #1 \\
     \end{pmatrix}
    \end{pmatrix}
}

\newcommand{\eg}{{\it e.g.}\xspace}
\newcommand{\etc}{{\it etc.}\xspace}
\newcommand{\ie}{{\it i.e.}\xspace}
\newcommand{\cf}{{\it cf.}\xspace}
\newcommand{\etal}{{\it et~al.}\xspace}
\newcommand{\ifandonlyif}{{\it iff}\xspace}
\newcommand{\resp}{{\it resp.}\xspace}

\newcommand{\bb}{\mathsf{bb}}
\newcommand{\be}{\mathsf{be}}
\newcommand{\ee}{\mathsf{ee}}

\newcommand{\trace}{\mathsf{Tr}_{\ExtensionField/\Field}}



\newcommand{\EqOrNot}{\mathrel{\overset{\mathrm{?}}{=}}}
%\DeclareUnicodeCharacter{\ast}{\op}

%\newcommand{\OP}{{\mathcal{O}\mathcal{P}}}
\newcommand{\op}{{\mathbin{\circledast}}}
\newcommand{\gtimes}{\mathbin{\odot}}
\newcommand{\gadd}{\mathbin{\oplus}}
\newcommand{\gminus}{\mathbin{\ominus}}
\newcommand{\gsum}{\mathop{\bigoplus}}
\newcommand{\kronecker}{\mathbin{\otimes}}
\newcommand{\notationmapsto}{\mathrel{\Rrightarrow}}
\newcommand{\notationmapsfrom}{\mathrel{\Lleftarrow}}

\newcommand{\gzero}{{O}}

\newcommand{\Group}{{\mathbb{G}}}
\newcommand{\Field}{{\mathbb{F}}}
\newcommand{\BaseField}{\mathbb{F}}
\newcommand{\ExtensionField}{\mathbb{E}}
\newcommand{\FieldMinusZO}{{\mathbb{F}\setminus\{0,1\}}}
\newcommand{\CodeSpace}{{\mathbb{C}}}
\newcommand{\Code}{{\mathcal{C}}}
\newcommand{\Zp}{{\mathbb{Z}}_p}
\newcommand{\TensorSpace}{\mathbb{T}}
\newcommand{\IndexSet}{{\mathcal{I}}}
\newcommand{\Root}{{\mathcal{R}}}
\newcommand{\Oracle}{{\mathcal{O}}}
\newcommand{\AFunc}{{\mathcal{F}}}


\newcommand{\ComSpace}{\mathbb{C}}
\newcommand{\HashSpace}{\mathbb{H}}
\newcommand{\PolySpace}{\mathbb{P}}
\newcommand{\FieldStar}{{\Field^*}}
\newcommand{\MSpace}{\mathbb{M}}
\newcommand{\QSpace}{\mathbb{Q}}

\newcommand{\ExpressionSpace}{\mathcal{E}} %only used in 4-expression

\newcommand{\security}{{1^{\lambda}}}

\newcommand{\slice}{{\langle:\rangle}}

\newcommand{\ChooseRandom}{\mathrel{\xleftarrow{\$}}}
\newcommand{\sample}{\mathrel{\xleftarrow{\$}}}
\newcommand{\Definition}{\mathrel{\overset{\text{def}}{=}}}

\newcommand{\Partition}{\xleftarrow{{\sf part}}}
\newcommand{\Round}{\xleftarrow{{\sf round}}}
\newcommand{\Flatten}{\xleftarrow{{\sf flat}}}
\newcommand{\Fold}{\xleftarrow{{\sf fold}}}
\newcommand{\Half}{\xleftarrow{{\sf half}}}
\newcommand{\Broadcast}{\xleftarrow{{\sf broadcast}}}
\newcommand{\Adversary}{\mathcal{A}}
\newcommand{\Prover}{\mathcal{P}}
\newcommand{\Verifier}{\mathcal{V}}
\newcommand{\Emulator}{\mathcal{E}}
\newcommand{\Simulator}{\mathcal{S}}
\newcommand{\Relation}{\mathcal{R}}
\newcommand{\Language}{\mathcal{L}}
\newcommand{\Extractor}{\mathcal{E}}
\newcommand{\Algorithm}{\mathcal{O}}
\newcommand{\Expect}{\mathbb{E}}
\newcommand{\Distinguisher}{\mathcal{D}}
\newcommand{\interact}{{\rightleftharpoons}}
\newcommand{\PVinteract}{{(\mathcal{P} \rightleftharpoons \mathcal{V})}}
\newcommand{\PtoV}{{\mathcal{P} \rightarrow \mathcal{V}}}
\newcommand{\VtoP}{{\mathcal{V} \rightarrow \mathcal{P}}}
\newcommand{\myParbox}[1]{\parbox{0.9\textwidth}{#1}}

\newcommand{\RelationExp}[3]{{\left\{~#1~:~#2~\left|~#3~\right.\right\}}}

%\newcommand{\flag}{{\sf flag}}
%\newcommand{\OpSet}{\mathcal{OP}}

\newcommand{\vectorize}[1]{{\sf vec}(#1)}
\newcommand{\Enc}{{\sf Enc}} %encoding
\newcommand{\diag}{{\sf diag}}
\newcommand{\tr}{{\sf tr}}
\newcommand{\rank}{{\sf rank}}
\newcommand{\col}{{\sf Col}}
\newcommand{\im}{{\sf Im}}
% \newcommand{\det}{{\sf det}}
\newcommand{\trans}{{\sf T}}

\newcommand{\Merkle}{{\sf Merkle}}


\newcommand{\trapdoor}{{\sf trap}} % only in appx-preliminary

\newcommand{\muls}{\mathbb{M}}
\newcommand{\hashes}{\mathsf{H}}
\newcommand{\mpath}{\mathsf{MP}}

\newcommand{\pccost}{{\sf pc}}
\newcommand{\zkcost}{{\sf zk}}
\newcommand{\exps}{\mathbb{E}}
\newcommand{\pairing}{\mathbb{P}} % only in table-evalyn-vs-zkml, table-appx-comp and table-appx-pcs

%\newcommand{\st}{\textrm{s.t.}}

\newcommand{\TRUE}{{\sf true}}
\newcommand{\FALSE}{{\sf false}}
\newcommand{\Accept}{{\sf Accept}}

\newcommand{\pc}{{\sf pc}}
\newcommand{\Polycom}{{\sf pc}}
%only keep \Polycom?

\newcommand{\io}{\bm{io}}
\newcommand{\vk}{{\sf vk}}
\newcommand{\pk}{{\sf pk}}

\newcommand{\tensorspan}[1]{ {\sf ts}(#1)}
\newcommand{\FS}{{\sf FS}}
\newcommand{\Setup}{{\sf Setup}}
\newcommand{\Open}{{\sf Open}}
\newcommand{\Query}{{\sf Query}}
\newcommand{\Prove}{{\sf Prove}}
\newcommand{\Verify}{{\sf Verify}}
\newcommand{\Eval}{{\sf EvalProof}}
\newcommand{\Com}{{\sf Comm}}
\newcommand{\Hash}{{\sf Hash}} %only in appx-preliminary
\newcommand{\Poly}{{\sf poly}}
\newcommand{\Negl}{{\sf negl}}
\newcommand{\pp}{{\sf pp}}
\newcommand{\rand}{{\sf rand}}
\newcommand{\bit}{{\sf bit}}

\newcommand{\suchthat}{\textrm{s.t.}}

%\newcommand{\params}{{\sf params}}
%\newcommand{\ID}{{\sf ID}}
%\newcommand{\len}{{\sf len}}
%\newcommand{\thread}{{\sf thread}}
%\newcommand{\Dec}{{\sf Dec}}
%\newcommand{\VerifyEval}{{\sf VerifyEval}}
%\newcommand{\VerifyPoly}{{\sf VerifyPoly}}
%\newcommand{\VerifyQuery}{{\sf VerifyQuery}}
%\newcommand{\VerifyCom}{{\sf VerifyCom}}
%\newcommand{\TP}{{\sf TP}}
%\newcommand{\TV}{{\sf TV}}
%\newcommand{\Solve}{{\sf Solve}}
%\newcommand{\range}{{\sf range}}
%\newcommand{\Coin}{{\sf Coin}}
%\newcommand{\HCom}{{\sf HComm}}
%\newcommand{\Time}{{\sf Time}}
%\newcommand{\Dist}{{\sf Dist}}
%\newcommand{\vp}{{\sf vp}}

\newcommand{\baseline}{{\#[{\sf baseline}]}}
\newcommand{\baselineAlg}[1]{{\#[{\sf baseline, alg{#1}}]}}

\newcommand{\mcircle}[1]{\mathop{\text{\textcircled{#1}}}}

% \makeatletter
% \newcommand\xleftrightarrow[2][]{%
% 	\ext@arrow 9999{\longleftrightarrowfill@}{#1}{#2}}
% \newcommand\longleftrightarrowfill@{%
% 	\arrowfill@\leftarrow\relbar\rightarrow}
% \makeatother

\newcommand{\customarrow}[1]{\xrightarrow{\makebox[#1]{}}}

\newtheorem{assumption}{Assumption}
\newtheorem{myclaim}{Claim}
\newtheorem{fact}{Fact}

\newtcolorbox{compactbox}[1][]{%
	colback=gray!5!white, % Background color
	colframe=black,		% Border color
	rounded corners,	% Rounded corners
	boxrule=0.5pt,		% Border thickness
	top=2pt,			% Top padding reduced
	bottom=2pt,		% Bottom padding reduced
	left=2pt,			% Left padding reduced
	right=2pt,		% Right padding reduced
	before skip=5pt,	% Space before the box reduced
	after skip=5pt,		% Space after the box reduced
	fonttitle=\bfseries,	% Title font bold
	title={\textbf{#1}},	% Use \textbf to make sure the title is bold
	attach title to upper, % Attach title inline
	% boxed title style={%
	%colback=gray!20!white, % Slightly darker background for title
	%colframe=black,	% Title frame color
	%top=-2pt,				% Title position adjustment
	%bottom=0pt,
	%left=4pt,
	%right=4pt,
	% },
	coltitle=black,		% Title text color
 % #1						% Additional user-defined options
}

\newtcolorbox{coloredalgo}[1][]{
	colback=gray!5!white, % Background color
	colframe=black, % Frame color
	boxrule=0.5pt, % Frame thickness
	#1 % Allows for further customizations passed at the environment call
}

\definecolor{mydarkblue}{HTML}{16468d}
\definecolor{mydarkgreen}{HTML}{226e4b}
\definecolor{myblue}{HTML}{bce2f5}
\definecolor{mygreen}{HTML}{bbead5}
\definecolor{myyellow}{HTML}{fff1dd}
\definecolor{mygrey}{HTML}{d9d9d9}

%\makeatletter
%\renewcommand\@cite[1]%{\leavevmode\unskip[\nobreakspace#1]\unskip\linebreak[0]}
%\makeatother

\newmdenv[
	backgroundcolor=gray!5,	% light gray background
	linecolor=white,		 % no border color
	innerleftmargin=0pt,	 % adjust inner margins
	innerrightmargin=0pt,
	innertopmargin=0pt,
	innerbottommargin=0pt,
	skipabove=5pt,			 % space above the box
	skipbelow=5pt,			% space below the box
        outermargin=0pt,
]{algorithmbackground}


\SetKwInput{KwPubIn}{Public input}
\SetKwInput{KwPriIn}{Private input}

\newlist{tightitemize}{itemize}{1}
\setlist[tightitemize]{label=\textbullet, leftmargin=*, topsep=0pt, partopsep=0pt, itemsep=0pt, parsep=0pt}

\begin{document}

\title{Mikros: SNARKs over Small Prime Fields}
% \subtitle{A}foldinverse
%Efficient zkSNARK for Matrix Expression Through Proof Composition 
% \title{Revisiting Bulletproofs: Surpassing Sumcheck in the zk-SNARK Campaign}
%Matilda: Matrix-Integrated Layered-Data Assurance/Authentication
%Miranda: Matrix-Integrated Recursive Non-interactive Atomic Derivation Assurance

%Sumcheck Replacement, Matrix Expressions, Commitment Reduction

%\titlerunning{}
\author{Paper \#01 \institute{\,} }
% \author{\, \institute{\,}}
% \author{Mingshu Cong
% \inst{\,\,*}
% \inst{1} \orcidID{0000-0001-5858-3411}
% \and
% Sherman S. M. Chow
% \inst{2}
% \orcidID{0000-0002-0629-6792}
% \and
% Tsz Hon Yuen
% \inst{3}
% \orcidID{0000-0002-0629-6792}
% \and
% Siu Ming Yiu
% \thanks{Siu Ming Yiu is the corresponding author.}
% \inst{1} \orcidID{0000-0002-3975-8500}}

% \authorrunning{M. Cong et al.}
% \institute{Department of Computer Science, The University of Hong Kong, Hong Kong SAR, China \\
% \email{\{mscong,smyiu\}@cs.hku.hk}
% \institute{Department of Computer Science, The Chinese University of Hong Kong, Hong Kong SAR, China \\
% \email{smchow@ie.cuhk.edu.hk}
% \institute{Faculty of Information Technology, Monash University, Melbourne, Australia\\
% \email{John.tszhonyuen@monash.edu}
% }
%
\maketitle

\begin{abstract}
Prover memory is a central bottleneck in modern succinct non-interactive
arguments of knowledge (SNARKs).  We construct SNARKs for arithmetic
circuits defined directly over a small prime field
\(\Field=\mathbb F_p\).  Keeping witnesses and intermediate protocol
states in the base field reduces the concrete memory cost per field
element, but also creates a soundness challenge: a single random
evaluation over \(\mathbb F_p\) provides only small-field security.

We address this problem with a tensor-structured multiplicative packing
of the LiteBullet inner-product argument.  For
\(h=\lceil\log_2\kappa\rceil\), Packed LiteBullet folds \(\kappa\)
same-length inner-product claims using \(\kappa\) independent challenges
in \(\mathbb F_p^h\times\mathbb F_p^\times\).  For
\(p-1>2(h+1)\), its round error is at most
\((2(h+1)/(p-1))^\kappa\), while all arithmetic remains in the base
field.
The folding matrices are generated on demand from these challenge
points.
For aggregate initial vector length \(N\) across the \(\kappa\)
channels, each channel has length \(m=N/\kappa\) and the folding depth
is \(d=\log_2 m\).  The prover uses
\(\Theta(N+\kappa^2)\) field elements of working memory and
\[
    O\left(
      \kappa N+
      \kappa^3d
    \right)
\]
field operations.  The packed claims may originate from unrelated
circuits with common padded dimensions, or from repeated copies of one
claim used for soundness amplification.

We combine this folding argument with a polynomial commitment layer
constructed from rank-one quasi-Abelian codes and Walsh--Hadamard
transforms, also directly over \(\mathbb F_p\).  These techniques
underlie \methodname, a SNARK framework designed for computations native
to small prime fields.
\end{abstract}

\section{Technical Overview}

We work throughout over a small odd prime field
\(\Field=\mathbb F_p\).  The circuit, witness, folding protocol, and
commitment scheme all use this same base field.  A computation with
\(N\) stored values therefore occupies
\(\Theta(N\log_2p)\) bits; when \(p<2^8\), each value fits in one byte.
This representation is particularly natural for computations already
arithmetized over a small field, including low-precision and quantized
workloads.

The main obstacle is soundness rather than arithmetic expressiveness.
After clearing its \(z\)-denominator, a nonzero round polynomial of
degree \(D\) vanishes at a uniform point of
\(\mathbb F_p^h\times\mathbb F_p^\times\) with probability at most
\(D/(p-1)\), which is too large when \(p\) is small.  Our solution is
to evaluate the folding relation at several independent base-field
points and retain the resulting claims in packed form.  This amplifies
soundness without changing the ambient field.

The construction has two algebraic layers.  Multiplicative Packed
LiteBullet compresses \(\kappa\) inner-product claims while preserving
\(\kappa\) independent challenges.  A code-based polynomial commitment
then binds the folded vector states and supports the terminal openings.
Both layers use tensor structure to avoid materializing their large
evaluation matrices.

For lookup constraints, Packed LogUp reshapes the query into
\(\kappa\) short columns, each with fewer than \(p\) entries and its own
multiplicity vector.  It combines the resulting rational identities
with a tensor challenge and evaluates the packed identity at
\(\kappa\) independent extension-field points.  These challenge pairs
create \(\kappa\) child columns whose aggregate length is the original
query length.  Subsequent rounds keep \(\kappa\) columns while halving
their length.  A packed reciprocal relation binds each child entry
directly to one row of the original query matrix, eliminating the
\(\kappa\)-fold inverse trace.

For permutation constraints, Packed Permutation replaces the usual
grand product by random base-field polynomials.  It first represents
each column by its first \(p\) power sums, then uses independent random
coefficient matrices to compress any nonzero moment residue into every
soundness lane.  The power tables are virtual and are checked through
degree-two Hadamard relations, so only the original columns are
committed.  This removes both the bit-decomposition witness and
extension-field arithmetic.  Exactness requires each packed column to
contain fewer than \(p\) entries, which is already the intended
small-channel regime.

\subsection{Computation over a Small Field}
\label{sec:small-field}

An arithmetic circuit is represented directly by the R1CS relation
\[
    (\bm A_L\bm w)\circ(\bm A_R\bm w)
       =\bm A_O\bm w,
\]
where
\(\bm A_L,\bm A_R,\bm A_O\in\Field^{m\times n}\) and
\(\bm w\in\Field^n\).  All equalities in the remainder of the paper are
over \(\Field=\mathbb F_p\) unless stated otherwise.

Small-field arithmetic admits a compact implementation.  For
\(a,b\in\mathbb F_p\), let \(s\leftarrow a+b\) in an integer type wide
enough to hold values below \(2p\).  Then
\[
    a+b\bmod p
    =
    \begin{cases}
        s-p, & s\ge p,\\
        s,   & s<p.
    \end{cases}
\]
Thus addition requires one integer addition followed by a conditional
subtraction.  For multiplication and inversion, we may precompute
\[
    \mathsf{Mul}[a,b]:=ab\bmod p,
    \qquad
    \mathsf{Inv}[a]:=a^{-1}\bmod p
    \quad(a\ne0).
\]
When \(p<2^8\), these tables require approximately \(p^2+p\) bytes in
total.  Vector and matrix operations then reduce to independent
small-integer operations and table lookups.

Direct small-field execution saves memory per element, but it also
restricts the verifier's challenge space.  One possible remedy is to
pack several base-field coordinates into one extension-field element.
We first record this standard algebraic reduction and clarify its cost.

\subsection{Packed Inner-Product Proof over an Extension Field}
\label{sec:extension-inner-product}

Consider a claimed inner product over \(\Field=\mathbb F_p\):
\[
    c\stackrel{?}{=}\langle\bm a,\bm b\rangle,
    \qquad
    \bm a,\bm b\in\Field^N.
\]
Let \(\kappa\mid N\), set \(m:=N/\kappa\), and partition each vector
into \(\kappa\) length-\(m\) blocks:
\[
    \bm a=
      \concat{\bm a^{(1)},\ldots,\bm a^{(\kappa)}},
    \qquad
    \bm b=
      \concat{\bm b^{(1)},\ldots,\bm b^{(\kappa)}}.
\]
The original relation is therefore
\begin{equation}
    c\stackrel{?}{=}
      \sum_{i=1}^{\kappa}
      \left\langle\bm a^{(i)},\bm b^{(i)}\right\rangle.
\label{eq:extension-packed-base-relation}
\end{equation}

\paragraph{The power basis and its constant-term dual.}
Let \(f(X)\in\Field[X]\) be a monic irreducible polynomial of degree
\(\kappa\), and define
\[
    \overline{\Field}
      :=\Field[X]/(f(X))
      \cong\mathbb F_{p^\kappa},
    \qquad
    \alpha:=X\bmod f(X).
\]
The tuple
\begin{equation}
    \mathcal B
      :=(1,\alpha,\ldots,\alpha^{\kappa-1})
\label{eq:extension-power-basis}
\end{equation}
is the power basis of \(\overline{\Field}\) over \(\Field\).  Calling
\(\mathcal B\) a basis means that every
\(\overline u\in\overline{\Field}\) has a unique representation
\[
    \overline u
      =u_0+u_1\alpha+\cdots+u_{\kappa-1}\alpha^{\kappa-1},
    \qquad
    u_0,\ldots,u_{\kappa-1}\in\Field.
\]
We write
\[
    \mathsf{ct}(\overline u)
      :=[\alpha^0]\overline u
      :=u_0
\]
for the \(\Field\)-linear map that extracts the constant coordinate in
this basis.

The constant-term dual basis
\(\mathcal B^\star=(\gamma_0,\ldots,\gamma_{\kappa-1})\) is defined by
\begin{equation}
    \mathsf{ct}(\alpha^i\gamma_j)=\delta_{i,j},
    \qquad
    0\leq i,j<\kappa.
\label{eq:extension-constant-dual-basis}
\end{equation}
It exists uniquely because the bilinear form
\((u,v)\mapsto\mathsf{ct}(uv)\) is non-degenerate: for every
\(u\ne0\), choosing \(v=u^{-1}\) gives \(\mathsf{ct}(uv)=1\).
The basis can be precomputed for any irreducible \(f\).  Define
\begin{equation}
    (\bm G)_{i,r}
      :=\mathsf{ct}\!\left(\alpha^{i+r}\right),
    \qquad
    0\leq i,r<\kappa.
\label{eq:extension-constant-gram-matrix}
\end{equation}
where multiplication is performed in \(\overline{\Field}\), and hence
is reduced modulo \(f\).  Then \(\bm G\) is invertible, and the
\(j\)-th dual element is
\begin{equation}
    \gamma_j
      :=\sum_{r=0}^{\kappa-1}
        (\bm G^{-1})_{r,j}\alpha^r.
\label{eq:extension-constant-dual-construction}
\end{equation}

\paragraph{A sparse dual basis.}
The basis conversion becomes particularly inexpensive when an
irreducible trinomial of the form
\begin{equation}
    f(X)=X^\kappa+\rho X+\sigma,
    \qquad \sigma\ne0,
\label{eq:extension-sparse-modulus}
\end{equation}
is available.  Since \(\alpha^\kappa=-\rho\alpha-\sigma\), the dual
basis has the explicit form
\begin{equation}
    \boxed{
    \gamma_0=1,
    \qquad
    \gamma_j=-\sigma^{-1}\alpha^{\kappa-j},
    \quad j=1,\ldots,\kappa-1.}
\label{eq:extension-sparse-dual-basis}
\end{equation}
The conditions involving \(j=0\) follow immediately from
\(\gamma_0=1\).  For \(j\geq1\), if \(i<j\), then
\(\alpha^i\gamma_j\) has degree strictly between zero and \(\kappa\).
If \(i=j\), then
\[
    \alpha^i\gamma_i
      =-\sigma^{-1}\alpha^\kappa
      =1+\rho\sigma^{-1}\alpha.
\]
Finally, if \(i>j\), writing \(r=i-j\geq1\) gives
\[
    \alpha^i\gamma_j
      =-\sigma^{-1}\alpha^{\kappa+r}
      =\alpha^r+\rho\sigma^{-1}\alpha^{r+1},
\]
which again has no constant term.  This proves
\eqref{eq:extension-constant-dual-basis} without requiring the cross
terms themselves to vanish.

\paragraph{Packing the inner product.}
Pack the blocks coordinate-wise using the power basis in the first
operand and its constant-term dual in the second:
\begin{equation}
\begin{aligned}
    \overline{\bm a}
      &:=\sum_{i=0}^{\kappa-1}
        \alpha^i\bm a^{(i+1)},\\
    \overline{\bm b}
      &:=\sum_{j=0}^{\kappa-1}
        \gamma_j\bm b^{(j+1)}
       \in\overline{\Field}^{m}.
\end{aligned}
\label{eq:extension-inner-product-packing}
\end{equation}
For the sparse dual basis in
\eqref{eq:extension-sparse-dual-basis}, the second equation is simply
\[
    \overline{\bm b}
      =\bm b^{(1)}
       -\sigma^{-1}
        \sum_{j=1}^{\kappa-1}
        \alpha^{\kappa-j}\bm b^{(j+1)};
\]
that is, its nonconstant coefficients are reversed and multiplied by
the single base-field scalar \(-\sigma^{-1}\).

Let
\(\overline c:=\langle\overline{\bm a},\overline{\bm b}\rangle\).
By bilinearity and~\eqref{eq:extension-constant-dual-basis},
\begin{equation}
\begin{aligned}
    \mathsf{ct}(\overline c)
      &=\sum_{i,j=0}^{\kappa-1}
        \left\langle
          \bm a^{(i+1)},\bm b^{(j+1)}
        \right\rangle
        \mathsf{ct}(\alpha^i\gamma_j)\\
      &=\sum_{i=1}^{\kappa}
        \left\langle\bm a^{(i)},\bm b^{(i)}\right\rangle
       =\langle\bm a,\bm b\rangle.
\end{aligned}
\label{eq:extension-inner-product-constant-term}
\end{equation}
Consequently, an extension-field proof of
\[
    \overline c
      \stackrel{?}{=}
      \langle\overline{\bm a},\overline{\bm b}\rangle
\]
together with the base-field linear check
\[
    c\stackrel{?}{=}\mathsf{ct}(\overline c)
\]
proves the original relation, provided that the packed vectors are
linked to the committed base-field vectors through
\eqref{eq:extension-inner-product-packing}.  This proves the aggregate
claim in~\eqref{eq:extension-packed-base-relation}; it does not by
itself certify separately supplied values
\(c^{(i)}=\langle\bm a^{(i)},\bm b^{(i)}\rangle\).

For our degree-\(16\) instantiation over \(\mathbb F_{251}\), we use
\begin{equation}
    \overline{\Field}
      :=\mathbb F_{251}[X]/(X^{16}+3X+1).
\label{eq:extension-concrete-field}
\end{equation}
The modulus is irreducible: Rabin's criterion gives
\[
    X^{251^{16}}\equiv X\pmod{f(X)},
    \qquad
    \gcd\!\left(f(X),X^{251^8}-X\right)=1.
\]
Writing \(\alpha:=X\bmod f(X)\), we have
\[
    \alpha^{16}=-3\alpha-1.
\]
The power basis and its constant-term dual are therefore
\begin{equation}
\begin{aligned}
    \mathcal B
      &=\left(1,\alpha,\alpha^2,\ldots,\alpha^{15}\right),\\
    \mathcal B^\star
      &=\left(1,-\alpha^{15},-\alpha^{14},\ldots,-\alpha\right).
\end{aligned}
\label{eq:extension-concrete-dual-basis}
\end{equation}
To see the duality explicitly, set
\(\gamma_0:=1\) and
\(\gamma_j:=-\alpha^{16-j}\) for \(1\leq j\leq15\).  For
\(j\geq1\),
\begin{equation}
    \alpha^i\gamma_j
      =
      \begin{cases}
        -\alpha^{16+i-j},
          & i<j,\\
        1+3\alpha,
          & i=j,\\
        \alpha^{i-j}+3\alpha^{i-j+1},
          & i>j.
      \end{cases}
\label{eq:extension-concrete-dual-cases}
\end{equation}
In the first and third cases, every exponent lies between \(1\) and
\(15\), so the constant coordinate is zero; in the middle case, it is
one.  The case \(j=0\) follows directly from \(\gamma_0=1\).  Hence
\[
    \mathsf{ct}(\alpha^i\gamma_j)=\delta_{i,j}
    \qquad(0\leq i,j\leq15).
\]
The concrete packing is consequently
\begin{equation}
\begin{aligned}
    \overline{\bm a}
      &=\sum_{i=0}^{15}\alpha^i\bm a^{(i+1)},\\
    \overline{\bm b}
      &=\bm b^{(1)}-
        \sum_{j=1}^{15}\alpha^{16-j}\bm b^{(j+1)}.
\end{aligned}
\label{eq:extension-concrete-packing}
\end{equation}
Thus the second packing requires only reversing and negating its
nonconstant base-field coordinates.  Moreover,
\[
    \lvert\overline{\Field}\rvert
      =251^{16}\approx2^{127.54},
\]
so one extension-field challenge provides an approximately
\(128\)-bit challenge space.

This reduction replaces a length-\(N\) base-field relation by a
length-\(m=N/\kappa\) extension-field relation.  Thus it uses
\(\kappa\) times fewer extension-field coordinates than naively
embedding all \(N\) entries into \(\overline{\Field}\).  This is not,
however, a \(\kappa\)-fold reduction in base-field work or memory.  If
one multiplication in \(\mathbb F_{p^\kappa}\) costs
\(\mathsf{M}(\kappa)\) base-field operations, the inner product costs
\[
    O\!\left(\frac{N}{\kappa}\mathsf{M}(\kappa)\right)
\]
base-field operations.  With the sparse dual basis above, packing costs
only \(O(N)\) base-field operations.  With schoolbook extension
multiplication, the inner product costs \(O(N\kappa)\), while fast
polynomial multiplication reduces it to \(\widetilde O(N)\).
Moreover, \(m\) extension elements contain exactly
\(\kappa m=N\) base-field coefficients, so materializing the packed
vectors does not asymptotically reduce their bit-level storage.

Our construction instead keeps the folding protocol over the base
field and obtains soundness from multiple small-field challenges.  It
follows the multiple-challenge principle of packed
sumcheck~\cite{iacr/WeiWXZDWZ25}, but replaces its univariate Lagrange
encoding by a tensor-structured multiplicative Laurent encoding over
\(\mathbb F_p\).

\subsection{Multiplicative Packed LiteBullet}
\subsection{Multiplicative Packed LiteBullet}
\label{sec:packed-litebullet}

\paragraph{State.}
Consider \(\kappa\) claimed inner products over
\(\Field=\mathbb F_p\):
\[
    c^{(i)}\stackrel{?}{=}\langle\bm a^{(i)},\bm b^{(i)}\rangle,
    \qquad i=1,\ldots,\kappa,
\]
where all vectors have the same power-of-two length \(m\).
At the initial round, we write
\[
    N:=\kappa m
    \qquad\text{and hence}\qquad
    m=\frac{N}{\kappa}
\]
for the aggregate vector length across all \(\kappa\) channels.
We assume that \(p\) is odd, so powers of two are invertible in
\(\Field\).
Set
\[
    h:=\left\lceil\log_2\kappa\right\rceil,
    \qquad
    d:=\log_2 m.
\]
For soundness amplification of one claim, the \(\kappa\) initial
claims are identical copies of that claim.
The construction uses an address space of size \(2^h\), but only
the \(\kappa\) occupied addresses correspond to actual claims; no dummy
claims are required.  We place the full vectors in claim-indexed
columns:
\[
    \bm A \coloneqq \concat{\bm{a}^{(1)}, \ldots, \bm{a}^{(\kappa)}},
    \quad
    \bm B \coloneqq \concat{\bm{b}^{(1)}, \ldots, \bm{b}^{(\kappa)}}
      \in\Field^{m \times \kappa},
    \qquad
    \bm c \coloneqq
      (c^{(1)}, \ldots, c^{(\kappa)})^\top
      \in\Field^\kappa.
\]
The batch relation is therefore
\begin{equation}
    \bm c
      \stackrel{?}{=}
      \diag(\bm A^\top\bm B).
\label{eq:packed-parent-matrix-relation}
\end{equation}

We then split the matrices by rows:
\[
    \bm A=
    \begin{pmatrix}
    \bm{A}_U\\
    \bm{A}_D
    \end{pmatrix}
    ,\quad
    \bm B=
    \begin{pmatrix}
    \bm{B}_U\\
    \bm{B}_D
    \end{pmatrix},
    \qquad
    \bm A_U,\bm A_D,\bm B_U,\bm B_D
      \in\Field^{m/2 \times \kappa}.
\]
Thus column \(i\) of \(\bm A_U\) is \(\bm a_U^{(i)}\), and similarly
for the other three blocks.  The \(i\)-th claim is equivalently
\begin{equation}
    c^{(i)}
      \stackrel{?}{=}
      \langle\bm a_{U}^{(i)},\bm b_{U}^{(i)}\rangle
      +
      \langle\bm a_{D}^{(i)} ,\bm b_{D}^{(i)}\rangle.
\label{eq:packed-parent-relation}
\end{equation}

\paragraph{Folding the claim columns and their \(U/D\) halves.}
The vector \(\bm y\in\Field^h\) folds the claim columns,
while an independent scalar \(z\in\Field^\times\) folds their \(U/D\)
halves.  Assign the \(\kappa\) claims distinct addresses
\[
    \bm s^{(1)},\ldots,\bm s^{(\kappa)}
      \in\{0,1\}^{h}.
\]
For \(\bm s=(s_1,\ldots,s_h)\), write
\(\bm y^{\bm s}:=\prod_{t=1}^{h}y_t^{s_t}\), and define
\begin{equation}
    \bm\eta_{\bm y}
      :=
      \left(
        \bm y^{\bm s^{(1)}},\ldots,
        \bm y^{\bm s^{(\kappa)}}
      \right)^\top
      \in\Field^\kappa.
\label{eq:wht-fold-definition}
\end{equation}
When \(\kappa=2^h\), this is the usual tensor vector
\(\bigotimes_{t=1}^{h}(1,y_t)\), up to the ordering of addresses.
For general \(\kappa\), the unused addresses simply do not appear in
\(\bm\eta_{\bm y}\); they create neither dummy channels nor additional
challenges.

The scalar \(z\) folds the two row blocks, using reciprocal weights
in the two operands.  This gives the length-\(m/2\) vectors
\begin{equation}
\begin{aligned}
    \bm a_{\bm y,z}
      &:=
      \left(
        \bm A_U+
        z\bm A_D
      \right) \bm{\eta}_{\bm y},\\
    \bm b_{\bm y,z}
      &:=
      \left(
        \bm B_U+
        z^{-1} \bm B_D
      \right) \bm{\eta}_{\bm y}.
\end{aligned}
\label{eq:wht-vector-fold}
\end{equation}

Define the honest round matrices
\begin{equation}
\begin{aligned}
    \bm C^\star&:=\bm A^\top\bm B,\\
    \bm U^\star&:=\bm A_D^\top\bm B_U,\\
    \bm D^\star&:=\bm A_U^\top\bm B_D.
\end{aligned}
\label{eq:wht-honest-round-matrices}
\end{equation}
Thus
\(\bm C^\star,\bm U^\star,\bm D^\star
  \in\Field^{\kappa\times\kappa}\).
Their folded inner product is the Laurent polynomial
\begin{equation}
\begin{aligned}
    \phi(\bm y,z)
      &:=
      \left\langle
        \bm a_{\bm y,z},
        \bm b_{\bm y,z}
      \right\rangle\\
      &=
      \bm\eta_{\bm y}^{\top}
      \left(\bm A_U+z\bm A_D\right)^\top
      \left(\bm B_U+z^{-1}\bm B_D\right)
      \bm\eta_{\bm y}\\
      &=
      \bm\eta_{\bm y}^{\top}
      \left(
        \bm C^\star
        +z\bm U^\star
        +z^{-1}\bm D^\star
      \right)
      \bm\eta_{\bm y}.
\end{aligned}
\label{eq:wht-folded-inner-product}
\end{equation}
Expanding the last expression gives
\begin{equation}
    \phi(\bm y,z)
      =
      \sum_{r,s=1}^{\kappa}
      \bm y^{\bm s^{(r)}+\bm s^{(s)}}
      \left(
        C^\star_{r,s}
        +zU^\star_{r,s}
        +z^{-1}D^\star_{r,s}
      \right).
\label{eq:wht-folded-inner-product-expansion}
\end{equation}
Each \(2\bm s^{(i)}\) is a distinct corner of
\(\{0,1,2\}^{h}\), and
\[
    \bm s^{(r)}+\bm s^{(s)}=2\bm s^{(i)}
    \quad\Longrightarrow\quad
    r=s=i.
\]
Therefore the corresponding \(z^0\) coefficient isolates exactly the
\(i\)-th parent inner product:
\begin{equation}
    [\bm y^{2\bm s^{(i)}}z^0]\,
      \phi(\bm y,z)
      =
      C^\star_{i,i}
      =
      \left\langle\bm a^{(i)},\bm b^{(i)}\right\rangle
      =:c^{\star(i)}.
\label{eq:wht-interpolation}
\end{equation}


\paragraph{Round message and random fold.}
The prover sends three matrices purporting to equal the honest
matrices in \eqref{eq:wht-honest-round-matrices}:
\[
    \widehat{\bm C},\widehat{\bm U},\widehat{\bm D}
      \in\Field^{\kappa\times\kappa}.
\]
The verifier first checks that the diagonal of the claimed parent
matrix agrees with the current state:
\begin{equation}
    \diag(\widehat{\bm C})\EqOrNot\bm c.
\label{eq:wht-parent-check}
\end{equation}
Only after these matrices are fixed, the verifier samples
\[
    (\bm y^{(1)},z^{(1)}),\ldots,
    (\bm y^{(\kappa)},z^{(\kappa)})
      \ChooseRandom\Field^h\times\Field^\times.
\]
For brevity, set
\[
    \bm\eta^{(j)}:=\bm\eta_{\bm y^{(j)}},
    \qquad
    \bm Y:=
      \begin{pmatrix}
        \bm\eta^{(1)}&\cdots&\bm\eta^{(\kappa)}
      \end{pmatrix},
    \qquad
    \bm Z:=\diag\!\left(
      z^{(1)},\ldots,z^{(\kappa)}
    \right).
\]
The next-round witness matrices are
\begin{equation}
\begin{aligned}
    \bm A'
      &:=
      \bm A_U\bm Y+\bm A_D\bm Y\bm Z,\\
    \bm B'
      &:=
      \bm B_U\bm Y+\bm B_D\bm Y\bm Z^{-1},
\end{aligned}
\qquad
    \bm A',\bm B'\in\Field^{(m/2)\times\kappa}.
\label{eq:wht-random-folds}
\end{equation}
Equivalently, their \(j\)-th columns are
\begin{equation}
\begin{aligned}
    \bm a'^{(j)}
      &=
      \left(\bm A_U+z^{(j)}\bm A_D\right)\bm\eta^{(j)},\\
    \bm b'^{(j)}
      &=
      \left(\bm B_U+(z^{(j)})^{-1}\bm B_D\right)\bm\eta^{(j)}.
\end{aligned}
\label{eq:wht-child-vectors}
\end{equation}
The corresponding child claims are defined by
\begin{equation}
\begin{aligned}
    c'^{(j)}
      &:=
      (\bm\eta^{(j)})^\top
      \left(
        \widehat{\bm C}
        +z^{(j)}\widehat{\bm U}
        +(z^{(j)})^{-1}\widehat{\bm D}
      \right)
      \bm\eta^{(j)},\\
    \bm c'
      &:=
      \left(c'^{(1)},\ldots,c'^{(\kappa)}\right)^\top.
\end{aligned}
\label{eq:wht-state-update}
\end{equation}
Equivalently, all child claims can be written at once as
\[
    \bm c'
      =
      \diag\!\left(
        \bm Y^\top\widehat{\bm C}\bm Y
        +\bm Z\bm Y^\top\widehat{\bm U}\bm Y
        +\bm Z^{-1}\bm Y^\top\widehat{\bm D}\bm Y
      \right).
\]
In the Fiat--Shamir transform, the \(y_t\)-coordinates are sampled
from \(\Field\), while rejection sampling ensures \(z^{(j)}\ne0\).
The round, child, and coordinate indices are included as domain
separators.

\paragraph{Folding invariant.}
Let
\[
    \bm c^\star:=
      \diag(\bm C^\star)
      =
      \left(
        \langle\bm a^{(1)},\bm b^{(1)}\rangle,
        \ldots,
        \langle\bm a^{(\kappa)},\bm b^{(\kappa)}\rangle
      \right)^\top.
\]
For the honest round message
\[
    (\widehat{\bm C},\widehat{\bm U},\widehat{\bm D})
      =
    (\bm C^\star,\bm U^\star,\bm D^\star),
\]
the parent and child states satisfy
\begin{equation}
    \boxed{
    \diag(\bm C^\star)=\bm c^\star,
    \qquad
    c'^{(j)}
      =
      \left\langle\bm a'^{(j)},\bm b'^{(j)}\right\rangle
      \quad\text{for every }j.
    }
\label{eq:wht-folding-invariant}
\end{equation}
Indeed, by \eqref{eq:wht-child-vectors},
\[
\begin{aligned}
    \left\langle\bm a'^{(j)},\bm b'^{(j)}\right\rangle
      &=
      (\bm\eta^{(j)})^\top
      \left(\bm A_U+z^{(j)}\bm A_D\right)^\top
      \left(\bm B_U+(z^{(j)})^{-1}\bm B_D\right)
      \bm\eta^{(j)}\\
      &=
      (\bm\eta^{(j)})^\top
      \left(
        \bm C^\star
        +z^{(j)}\bm U^\star
        +(z^{(j)})^{-1}\bm D^\star
      \right)
      \bm\eta^{(j)}
      =c'^{(j)}.
\end{aligned}
\]
Thus one round maps \(\kappa\) inner products of length \(m\) to
\(\kappa\) inner products of length \(m/2\).

\paragraph{One-claim sanity check.}
When \(\kappa=1\), we have \(h=0\) and
\(\bm\eta_{\bm y}=1\).  The parent entry is
\(C^\star=\langle\bm a,\bm b\rangle\), while the child claim is
\[
    c'=C^\star+zU^\star+z^{-1}D^\star
      =
      \left\langle
        \bm a_U+z\bm a_D,\,
        \bm b_U+z^{-1}\bm b_D
      \right\rangle,
\]
which is the usual multiplicative LiteBullet fold.

\paragraph{Round soundness.}
Assume that the vectors are binding-linked to the committed witness
and that at least one parent claim is false, so
\(\bm c\ne\bm c^\star\).  Suppose nevertheless that
\eqref{eq:wht-parent-check} accepts.  Define
\[
    \bm E_C:=\widehat{\bm C}-\bm C^\star,
    \qquad
    \bm E_U:=\widehat{\bm U}-\bm U^\star,
    \qquad
    \bm E_D:=\widehat{\bm D}-\bm D^\star.
\]
Then
\[
    \diag(\bm E_C)=\bm c-\bm c^\star\ne\bm0.
\]
Consider the Laurent polynomial
\begin{equation}
    \Delta(\bm y,z)
      :=
      \bm\eta_{\bm y}^\top
      \left(\bm E_C+z\bm E_U+z^{-1}\bm E_D\right)
      \bm\eta_{\bm y}.
\label{eq:wht-error-polynomial}
\end{equation}
This polynomial is nonzero.  In fact, if the \(i\)-th parent claim is
false, the corner property in \eqref{eq:wht-interpolation} gives
\[
    [\bm y^{2\bm s^{(i)}}z^0]\,
      \Delta(\bm y,z)
      =
      (\bm E_C)_{i,i}
      =
      c^{(i)}-c^{\star(i)}
      \ne0.
\]
Moreover, \(z\Delta(\bm y,z)\) is an ordinary polynomial with
individual degree at most two in each variable and total degree at
most \(2(h+1)\).

The state update gives the exact child-error identity
\begin{equation}
    c'^{(j)}
      -
      \left\langle\bm a'^{(j)},\bm b'^{(j)}\right\rangle
      =
      \Delta(\bm y^{(j)},z^{(j)}).
\label{eq:wht-child-error-identity}
\end{equation}
Consequently, all false information disappears in this round only if
all \(\kappa\) independent challenges are roots of the same nonzero
polynomial.  Assuming \(p-1>2(h+1)\), Schwartz--Zippel over
\(\Field^h\times\Field^\times\) yields
\[
    \boxed{
    \epsilon_{\mathsf{round}}
      \leq
      \left(\frac{2(h+1)}{p-1}\right)^\kappa,
      \qquad
      h=\left\lceil\log_2\kappa\right\rceil.
    }
\]
For aggregate initial vector length \(N\), each channel has length
\(m=N/\kappa\).  A union bound over the \(d\) rounds gives
\[
    \epsilon_{\mathsf{fold}}
      \leq
      d
      \left(\frac{2(h+1)}{p-1}\right)^\kappa.
\]
The complete argument additionally includes the soundness errors of
the terminal commitment openings and linking checks.  The same proof
applies to Fiat--Shamir in the random-oracle model, provided that the
round message is absorbed before the domain-separated challenges are
derived.

\paragraph{Cost and space.}
The round message contains \(3\kappa^2\) field elements.  Because
\(\diag(\widehat{\bm C})=\bm c\) is fixed by the parent state, these
\(\kappa\) diagonal entries can be omitted, reducing the transmitted
message to
\[
    3\kappa^2-\kappa
\]
field elements per round.  Straightforward dense computation of
\(\bm C^\star,\bm U^\star,\bm D^\star\) costs
\(O(\kappa^2m)\) field operations.  Constructing
\(\bm A'\) and \(\bm B'\) also costs \(O(\kappa^2m)\), and evaluating
the \(\kappa\) quadratic forms in \eqref{eq:wht-state-update} costs
\(O(\kappa^3)\).  Hence one round costs
\[
    O\left(\kappa^2m+\kappa^3\right).
\]
Initially \(m=N/\kappa\).  Since the per-channel lengths halve
geometrically over the \(d\) rounds, the total
arithmetic cost is
\[
    O\left(
      \kappa N+\kappa^3d
    \right),
\]
and the total folding communication is
\[
    (3\kappa^2-\kappa)d
\]
field elements.

At vector length \(m\), the current witness matrices contain
\(2\kappa m\) field elements.  The three round matrices and the
folding matrix \(\bm Y\) require \(O(\kappa^2)\) additional elements;
\(\bm Z\) is diagonal and can be generated on demand.  An in-place
implementation therefore uses
\[
    2\kappa m+O(\kappa^2)
\]
field elements in a round, and its peak working memory is
\[
    \Theta\!\left(N+\kappa^2\right).
\]
For independently supplied claims, the \(2N\) input field elements
already dominate this memory bound whenever \(N=\Omega(\kappa^2)\).
These bounds count only the algebraic folding layer.  It remains to
bind the terminal folded values to the initially committed matrices.
We next reduce this linking relation to the MLE interface used by the
commitment layer below.

\paragraph{Terminal MLE linking.}
Write
\[
    (\bm A_0,\bm B_0,\bm c_0)
      :=(\bm A,\bm B,\bm c)
\]
for the initial state, and let \(\bm c_\ell\) denote the claim vector
after round \(\ell\).  For round \(\ell=1,\ldots,d\), let
\(\bm y_\ell^{(j)},z_\ell^{(j)}\) be the challenge point of channel \(j\),
and let \(\bm Y_\ell,\bm Z_\ell\) be the corresponding matrices defined as
in \eqref{eq:wht-random-folds}.  The folded states satisfy
\begin{equation}
\begin{aligned}
    \bm A_\ell
      &=
      (\bm A_{\ell-1})_U\bm Y_\ell
      +(\bm A_{\ell-1})_D\bm Y_\ell\bm Z_\ell,\\
    \bm B_\ell
      &=
      (\bm B_{\ell-1})_U\bm Y_\ell
      +(\bm B_{\ell-1})_D\bm Y_\ell\bm Z_\ell^{-1},
\end{aligned}
\qquad
    \bm A_\ell,\bm B_\ell
      \in\Field^{(m/2^\ell)\times\kappa}.
\label{eq:wht-round-indexed-folds}
\end{equation}
At \(\ell=d\), both matrices contain one row.  Denote their channel
entries by
\[
    \bm A_d=
      \begin{pmatrix}
        a_d^{(1)}&\cdots&a_d^{(\kappa)}
      \end{pmatrix},
    \qquad
    \bm B_d=
      \begin{pmatrix}
        b_d^{(1)}&\cdots&b_d^{(\kappa)}
      \end{pmatrix}.
\]
The terminal inner-product relation is therefore
\begin{equation}
    c_d^{(j)}
      \stackrel{?}{=}
      a_d^{(j)}b_d^{(j)},
    \qquad j=1,\ldots,\kappa.
\label{eq:wht-terminal-products}
\end{equation}

To make the dependence on the initial matrices explicit, index their
rows by \(\bm x=(x_1,\ldots,x_d)\in\{0,1\}^d\).  Unrolling
\eqref{eq:wht-round-indexed-folds} gives
\begin{equation}
\begin{aligned}
    a_d^{(j)}
      &=
      \sum_{\bm x\in\{0,1\}^d}
      \sum_{i=1}^{\kappa}
      A_{\bm x,i}\,
      \bm e_i^\top
      \left(
        \prod_{\ell=1}^{d}
          \bm Y_\ell\bm Z_\ell^{x_\ell}
      \right)
      \bm e_j,\\
    b_d^{(j)}
      &=
      \sum_{\bm x\in\{0,1\}^d}
      \sum_{i=1}^{\kappa}
      B_{\bm x,i}\,
      \bm e_i^\top
      \left(
        \prod_{\ell=1}^{d}
          \bm Y_\ell\bm Z_\ell^{-x_\ell}
      \right)
      \bm e_j.
\end{aligned}
\label{eq:wht-terminal-linear-maps}
\end{equation}
Here the products are ordered from \(\ell=1\) to \(\ell=d\).  Thus every
terminal entry is a public, transcript-dependent linear functional of
one of the two initial matrices.  Notice that these coefficients are
generally not a tensor-span vector: the dense matrices \(\bm Y_\ell\)
are interleaved with the row-folding matrices \(\bm Z_\ell^{\pm1}\).
Consequently, the terminal entries are not themselves MLE evaluations
of \(\bm A_0\) or \(\bm B_0\).

\paragraph{Folding the terminal rows back to the committed inputs.}
We link the terminal rows to \(\bm A\) and \(\bm B\) directly by
folding; no general circuit-evaluation reduction is needed.  Let
\(\bm e_j\in\Field^\kappa\) denote the \(j\)-th standard basis vector,
whose only nonzero entry is a one in position \(j\).  Thus
\[
    \bm A_d\bm e_j=a_d^{(j)},
    \qquad
    \bm B_d\bm e_j=b_d^{(j)}.
\]
For terminal channel \(j\), initialize two public test matrices
\[
    \bm R_d^{A,(j)}:=\bm e_j^\top,
    \qquad
    \bm R_d^{B,(j)}:=\bm e_j^\top.
\]
For \(\ell=d,d-1,\ldots,1\), fold these test matrices backward through
the transpose of the corresponding forward map:
\begin{equation}
\begin{aligned}
    \bm R_{\ell-1}^{A,(j)}
      &:={}
      \begin{pmatrix}
        \bm R_\ell^{A,(j)}\bm Y_\ell^\top\\
        \bm R_\ell^{A,(j)}\bm Z_\ell\bm Y_\ell^\top
      \end{pmatrix},\\
    \bm R_{\ell-1}^{B,(j)}
      &:={}
      \begin{pmatrix}
        \bm R_\ell^{B,(j)}\bm Y_\ell^\top\\
        \bm R_\ell^{B,(j)}\bm Z_\ell^{-1}\bm Y_\ell^\top
      \end{pmatrix}.
\end{aligned}
\label{eq:wht-adjoint-folds}
\end{equation}
For example, the first backward fold is explicitly
\[
\begin{aligned}
    \bm R_{d-1}^{A,(j)}
      &=
      \begin{pmatrix}
        \bm e_j^\top\bm Y_d^\top\\
        \bm e_j^\top\bm Z_d\bm Y_d^\top
      \end{pmatrix},
    &
    \bm R_{d-1}^{B,(j)}
      &=
      \begin{pmatrix}
        \bm e_j^\top\bm Y_d^\top\\
        \bm e_j^\top\bm Z_d^{-1}\bm Y_d^\top
      \end{pmatrix}.
\end{aligned}
\]
After all \(d\) backward folds, the test matrices can be written
without recursion as
\begin{equation}
\begin{aligned}
    \bm R_0^{A,(j)}
      &=
      \left(
        \bm e_i^\top
        \left(
          \prod_{\ell=1}^{d}
            \bm Y_\ell\bm Z_\ell^{x_\ell}
        \right)
        \bm e_j
      \right)_{
        \bm x\in\{0,1\}^d,\,i\in[\kappa]
      },\\
    \bm R_0^{B,(j)}
      &=
      \left(
        \bm e_i^\top
        \left(
          \prod_{\ell=1}^{d}
            \bm Y_\ell\bm Z_\ell^{-x_\ell}
        \right)
        \bm e_j
      \right)_{
        \bm x\in\{0,1\}^d,\,i\in[\kappa]
      }.
\end{aligned}
\label{eq:wht-explicit-test-matrices}
\end{equation}
Here the outer indices mean that the displayed scalar is placed in row
\(\bm x\) and column \(i\).  Hence every entry of \(\bm R_0^{A,(j)}\)
or \(\bm R_0^{B,(j)}\) is exactly the coefficient with which the
corresponding entry of \(\bm A\) or \(\bm B\) contributes to terminal
entry \(j\).

Writing \(\langle\cdot,\cdot\rangle_F\) for the Frobenius inner
product, the forward and backward folds satisfy the exact invariants
\begin{equation}
\begin{aligned}
    \left\langle
      \bm A_\ell,\bm R_\ell^{A,(j)}
    \right\rangle_F
      &=
    \left\langle
      \bm A_{\ell-1},\bm R_{\ell-1}^{A,(j)}
    \right\rangle_F,\\
    \left\langle
      \bm B_\ell,\bm R_\ell^{B,(j)}
    \right\rangle_F
      &=
    \left\langle
      \bm B_{\ell-1},\bm R_{\ell-1}^{B,(j)}
    \right\rangle_F.
\end{aligned}
\label{eq:wht-adjoint-invariant}
\end{equation}
Consequently, the two terminal entries are linked to the initially
committed matrices by the dense linear claims
\begin{equation}
\begin{aligned}
    a_d^{(j)}
      &\stackrel{?}{=}
      \left\langle\bm A,\bm R_0^{A,(j)}\right\rangle_F,\\
    b_d^{(j)}
      &\stackrel{?}{=}
      \left\langle\bm B,\bm R_0^{B,(j)}\right\rangle_F.
\end{aligned}
\label{eq:wht-dense-terminal-links}
\end{equation}
The matrices \(\bm R_0^{A,(j)}\) and \(\bm R_0^{B,(j)}\) depend only
on the public folding transcript.  They need not be materialized: their
entries, or their folded values below, can be generated on demand from
\(\bm Y_1,\bm Z_1,\ldots,\bm Y_d,\bm Z_d\).

\paragraph{Folding a dense linear claim to a monomial-form evaluation.}
Consider either claim in \eqref{eq:wht-dense-terminal-links}.  Pad the
unused channel addresses with zeros and flatten the two matrices into
vectors \(\bm x,\bm r\in\Field^{2^{d+h}}\), so the current claim is
\[
    \alpha\stackrel{?}{=}\langle\bm x,\bm r\rangle.
\]
For this linking interface, we use the monomial-form multilinear
evaluation.  For \(\bm q=(q_1,\ldots,q_t)\), write
\begin{equation}
    \bm\mu_{\bm q}
      :=
      \bigotimes_{i=1}^{t}(1,q_i)
      \in\Field^{2^t}.
\label{eq:wht-monomial-tensor}
\end{equation}
Throughout this reduction, ``multilinear evaluation'' refers to the
inner product with \(\bm\mu_{\bm q}\) in
\eqref{eq:wht-monomial-tensor}.

In one folding step, split both vectors into their two halves.  The
prover sends the cross terms
\[
    \gamma_{UD}:=\langle\bm x_U,\bm r_D\rangle,
    \qquad
    \gamma_{DU}:=\langle\bm x_D,\bm r_U\rangle.
\]
After these values are fixed, the verifier samples
\(\rho\ChooseRandom\Field^\times\) and both parties set
\begin{equation}
\begin{aligned}
    \bm x'
      &:=\bm x_U+\rho\bm x_D,\\
    \bm r'
      &:=\bm r_U+\rho^{-1}\bm r_D,\\
    \alpha'
      &:=\alpha
        +\rho^{-1}\gamma_{UD}
        +\rho\gamma_{DU}.
\end{aligned}
\label{eq:wht-link-fold}
\end{equation}
The update preserves the inner-product relation:
\[
    \alpha'\stackrel{?}{=}\langle\bm x',\bm r'\rangle.
\]
For each terminal channel, the \(\bm A\)- and \(\bm B\)-linking
instances use the same challenge \(\rho\) in every step.  Repeating the
fold over the \(d\) row coordinates and the \(h\) channel coordinates
therefore leaves a common point
\[
    (\bm r_{\mathsf{row}}^{(j)},
      \bm r_{\mathsf{ch}}^{(j)})
      \in\Field^d\times\Field^h.
\]
Because the \(\bm x\)-side uses the weights \((1,\rho)\), its
final value is exactly an inner product with the monomial tensor in
\eqref{eq:wht-monomial-tensor}.  The two linking instances end with
\begin{equation}
\begin{aligned}
    v_A^{(j)}
      &\stackrel{?}{=}
      \bm\mu_{\bm r_{\mathsf{row}}^{(j)}}^\top
      \bm A
      \bm\mu_{\bm r_{\mathsf{ch}}^{(j)}},\\
    v_B^{(j)}
      &\stackrel{?}{=}
      \bm\mu_{\bm r_{\mathsf{row}}^{(j)}}^\top
      \bm B
      \bm\mu_{\bm r_{\mathsf{ch}}^{(j)}}.
\end{aligned}
\label{eq:wht-terminal-ab-mle-claims}
\end{equation}
The reciprocal fold of the public \(\bm r\)-side produces a known
scalar, and the verifier checks the product of this scalar and
\(v_A^{(j)}\), or \(v_B^{(j)}\), against the corresponding final folded
claim.

To combine the two openings, define
\[
    \bm V\in\Field^{m\times2^{h+1}}
\]
with columns indexed by
\((b,\bm s)\in\{0,1\}\times\{0,1\}^h\), and set
\begin{equation}
    V_{\bm x,(0,\bm s^{(i)})}:=A_{\bm x,i},
    \qquad
    V_{\bm x,(1,\bm s^{(i)})}:=B_{\bm x,i},
\label{eq:wht-combined-input}
\end{equation}
with all unused channel addresses set to zero.  After
\(v_A^{(j)}\) and \(v_B^{(j)}\) are fixed, the verifier samples
\(\beta^{(j)}\ChooseRandom\Field^\times\) and sets
\[
    v^{(j)}
      :=
      v_A^{(j)}+\beta^{(j)}v_B^{(j)}.
\]
The monomial tensor in the selector coordinate gives the single
opening claim
\begin{equation}
    v^{(j)}
      \stackrel{?}{=}
      \bm\mu_{\bm r_{\mathsf{row}}^{(j)}}^\top
      \bm V
      \left(
        \begin{pmatrix}1\\ \beta^{(j)}\end{pmatrix}
        \otimes
        \bm\mu_{\bm r_{\mathsf{ch}}^{(j)}}
      \right),
    \qquad j=1,\ldots,\kappa.
\label{eq:wht-terminal-mle-claims}
\end{equation}
Thus the two terminal rows reduce to exactly \(\kappa\) monomial-form
multilinear-evaluation claims on one committed matrix \(\bm V\).

\paragraph{Soundness of the terminal linking fold.}
Suppose a current inner-product claim has nonzero error \(\delta\), and
let \(\Delta_{UD}\) and \(\Delta_{DU}\) be the errors in the two cross
terms sent by the prover.  The child claim can become correct only if
\begin{equation}
    \rho\delta
      +\Delta_{UD}
      +\rho^2\Delta_{DU}
      =0.
\label{eq:wht-link-error-polynomial}
\end{equation}
This is a nonzero polynomial of degree at most two: if it were
identically zero, its constant and linear coefficients would imply
\(\Delta_{UD}=0\) and then \(\delta=0\).  Hence one fold has error at
most \(2/(p-1)\).  A union bound over the \(d+h\) folds, followed by
the degree-one \(\beta\)-combination of the final \(\bm A\)- and
\(\bm B\)-openings, gives
\begin{equation}
    \boxed{
    \epsilon_{\mathsf{link},1}
      \leq
      \frac{2(d+h)+1}{p-1}.
    }
\label{eq:wht-link-soundness}
\end{equation}
This is the soundness bound when every terminal entry is linked once;
the presence of \(\kappa\) distinct terminal entries alone does not
raise this bound to the \(\kappa\)-th power, because only one of those
entries may be false.  If the same linking claim is re-packed in
\(\kappa\) lanes at every fold and the resulting common nonzero error
polynomial is evaluated at \(\kappa\) independent challenges, then
\begin{equation}
    \epsilon_{\mathsf{link},\kappa}
      \leq
      (d+h)\left(\frac{2}{p-1}\right)^\kappa
      +\left(\frac{1}{p-1}\right)^\kappa
      \leq
      \left(\frac{2(d+h+1)}{p-1}\right)^\kappa.
\label{eq:wht-link-amplified-soundness}
\end{equation}
The soundness error of the final batched commitment openings is added
separately.

\subsection{Small Field Lookup Proof}
\label{sec:lookup}

Let
\[
    \bm v=(v_1,\ldots,v_N)^\top\in\Field^N,
    \qquad
    \bm t=(t_1,\ldots,t_M)^\top\in\Field^M
\]
be the lookup vector and the table, respectively.  For clarity, assume
that \(N\) and \(M\) are powers of two, and put
\[
    n:=\log_2N,
    \qquad
    d:=\log_2M.
\]
The general case is obtained by padding to the next power of two and
constraining every selected address to lie in the unpadded range.

\paragraph{Lookup as a virtual selection matrix.}
A lookup witness specifies, for every \(i\in[N]\), an address
\(\iota_i\in[M]\) such that \(v_i=t_{\iota_i}\).  Equivalently, it
specifies a Boolean selection matrix
\(\bm\Sigma\in\{0,1\}^{N\times M}\) having exactly one \(1\) in
each row, with
\begin{equation}
    \bm v=\bm\Sigma\bm t.
\label{eq:small-field-lookup-relation}
\end{equation}
We do not materialize or commit to \(\bm\Sigma\).  Instead, let
\[
    \bm B=(B_{i,j})
      \in\{0,1\}^{N\times d}
\]
contain the binary addresses: row \(i\) is the \(d\)-bit encoding of
\(\iota_i-1\).  The constraint
\begin{equation}
    \bm B\circ(\bm B-\bm1)=\bm0
\label{eq:small-field-lookup-booleanity}
\end{equation}
ensures that every entry of \(\bm B\) is Boolean and hence that every
row defines one standard-basis row of \(\bm\Sigma\).

\paragraph{Reducing the lookup to two inner products.}
After \(\bm v\), \(\bm t\), and \(\bm B\) are bound, the verifier
samples \(\bm x_l\in\Field^n\) and sets
\[
    \bm l
      :=\bigotimes_{j=1}^{n}(1,x_{l,j})
      \in\Field^N.
\]
The prover computes and commits to the pushforward
\begin{equation}
    \bm s:=\bm\Sigma^\top\bm l\in\Field^M.
\label{eq:small-field-lookup-pushforward}
\end{equation}
If \eqref{eq:small-field-lookup-relation} holds, then
\begin{equation}
    \boxed{
    \langle\bm l,\bm v\rangle
      =\langle\bm s,\bm t\rangle.}
\label{eq:small-field-lookup-table-check}
\end{equation}
It remains to prove that the committed vector \(\bm s\) is indeed the
pushforward in \eqref{eq:small-field-lookup-pushforward}.

After \(\bm s\) is bound, the verifier samples
\(\bm x_r\in\Field^d\) and defines the tensor vector
\begin{equation}
    \bm r
      :=\bigotimes_{j=1}^{d}(1,x_{r,j})
      \in\Field^M.
\label{eq:small-field-lookup-right-tensor}
\end{equation}
Writing
\(\widetilde{\bm r}:=\bm\Sigma\bm r\in\Field^N\), the adjoint
identity gives
\begin{equation}
    \boxed{
    \langle\bm s,\bm r\rangle
      =\langle\bm l,\widetilde{\bm r}\rangle.}
\label{eq:small-field-lookup-pushforward-check}
\end{equation}
Thus the lookup is reduced to
\eqref{eq:small-field-lookup-table-check} and
\eqref{eq:small-field-lookup-pushforward-check}, provided that
\(\widetilde{\bm r}\) is correctly derived from the committed address
bits.

\paragraph{Evaluating the virtual selection matrix.}
Let
\[
    \bm X:=\bm1_N\bm x_r^\top\in\Field^{N\times d},
    \qquad
    \bm R:=\bm1+\bm B\circ(\bm X-\bm1)
      \in\Field^{N\times d}.
\]
Entrywise,
\[
    R_{i,j}=1+B_{i,j}(x_{r,j}-1)
      =
      \begin{cases}
        1,&B_{i,j}=0,\\
        x_{r,j},&B_{i,j}=1.
      \end{cases}
\]
Consequently, row \(i\) selects the tensor entry addressed by its bit
string:
\begin{equation}
    \widetilde r_i
      =\prod_{j=1}^{d}R_{i,j}
      =\prod_{j=1}^{d}x_{r,j}^{B_{i,j}}
      =r_{\iota_i}.
\label{eq:small-field-lookup-row-product}
\end{equation}
Equivalently,
\[
    \widetilde{\bm r}
      =\mathsf{RowProd}(\bm R)
      =\bm\Sigma\bm r.
\]

\paragraph{Proving the row products.}
The prover commits to the prefix-product matrix
\[
    \bar{\bm R}\in\Field^{N\times(d-1)},
    \qquad
    \bar R_{i,j}:=\prod_{h=1}^{j}R_{i,h},
\]
and to \(\widetilde{\bm r}\in\Field^N\).  All row products are then
captured by the single Hadamard identity
\begin{equation}
    \boxed{
    \left(\bm1_N\,\middle|\,\bar{\bm R}\right)
      \circ\bm R
      =
    \left(\bar{\bm R}\,\middle|\,\widetilde{\bm r}\right).}
\label{eq:small-field-lookup-prefix-products}
\end{equation}
Both sides of \eqref{eq:small-field-lookup-prefix-products} have
dimension \(N\times d\).  The first column checks
\(\bar R_{i,1}=R_{i,1}\), every intermediate column checks the next
prefix product, and the last column checks
\(\widetilde r_i=\bar R_{i,d-1}R_{i,d}\).

The matrix \(\bm R\) need not be committed separately: it is an affine
function of the committed \(\bm B\) and the public challenge
\(\bm x_r\).  After commitments to \(\bm s\),
\(\bar{\bm R}\), and \(\widetilde{\bm r}\) are fixed, the verifier
checks
\eqref{eq:small-field-lookup-booleanity},
\eqref{eq:small-field-lookup-prefix-products},
\eqref{eq:small-field-lookup-table-check}, and
\eqref{eq:small-field-lookup-pushforward-check}.  These matrix and
inner-product claims can be accumulated into the Packed LiteBullet
interface of Section~\ref{sec:packed-litebullet}.

\paragraph{Soundness and cost.}
If \(\bm v\ne\bm\Sigma\bm t\), then
\[
    \langle\bm l,\bm v-\bm\Sigma\bm t\rangle
\]
is a nonzero multilinear polynomial in \(\bm x_l\) of degree at most
\(n\).  If \(\bm s\ne\bm\Sigma^\top\bm l\), then
\[
    \langle
      \bm s-\bm\Sigma^\top\bm l,
      \bm r
    \rangle
\]
is a nonzero multilinear polynomial in \(\bm x_r\) of degree at most
\(d\).  Hence the two tensor tests contribute at most
\begin{equation}
    \epsilon_{\mathsf{lookup}}
      \leq
      \min\left\{1,\frac{n+d}{p}\right\}
      +\epsilon_{\mathsf{row}}
      +\epsilon_{\mathsf{com}},
\label{eq:small-field-lookup-soundness}
\end{equation}
where \(\epsilon_{\mathsf{row}}\) is the error of the Booleanity and
prefix-product checks and \(\epsilon_{\mathsf{com}}\) collects the
binding and opening errors.  Repeating the complete tensor test with
independent Fiat--Shamir challenges and packing the repetitions into
independent channels amplifies the tensor-test error multiplicatively.

For one repetition, the protocol commits to \(Nd\) address bits,
\(N(d-1)\) prefix products, the length-\(N\) vector
\(\widetilde{\bm r}\), and the length-\(M\) pushforward \(\bm s\).
Thus its auxiliary witness size is \(O(N\log M+M)\) base-field
elements.  With \(\kappa\) parallel repetitions, \(\bm B\) is shared,
whereas the challenge-dependent prefix products, selected vectors, and
pushforwards occupy \(O(\kappa(N\log M+M))\) elements.  The advantage
is that every local algebraic relation has degree at most two, while
the two global fingerprints have degrees only \(\log_2N\) and
\(\log_2M\), rather than \(N\) or \(M\).

\subsection{Packed LogUp}
\label{sec:packed-logup}

Logarithmic-derivative lookup arguments reduce a lookup relation to an
equality between sums of reciprocals~\cite{iacr/Habock22a}.  Over a
small field, however, a residue can disappear modulo the
characteristic.  For example,
\[
    \underbrace{\frac{1}{X+v}+\cdots+\frac{1}{X+v}}
      _{p\text{ terms}}
    =0.
\]
We avoid this \emph{characteristic overflow} by placing fewer than
\(p\) query occurrences in each initial column.  The \(\kappa\) column
identities are first combined symbolically and then evaluated at
\(\kappa\) independent pairs of challenges.  Those evaluations form
\(\kappa\) new columns, which are recursively repacked exactly as in
the folding protocol above.

\paragraph{Packing the query into \(\kappa\) columns.}
Let
\[
    \bm v=(v_1,\ldots,v_n)^\top\in\Field^n,
    \qquad
    \bm t=(t_1,\ldots,t_T)^\top\in\Field^T
\]
be the query and table columns.  We assume that the table values are
distinct; duplicate table entries can be merged.  The goal is to prove
\[
    v_i\in\{t_1,\ldots,t_T\}
    \qquad\text{for all }i\in[n].
\]
Set
\[
    m:=\left\lceil\frac n\kappa\right\rceil,
    \qquad
    h:=\left\lceil\log_2\kappa\right\rceil,
\]
and assume
\begin{equation}
    m<p.
\label{eq:packed-logup-short-column-condition}
\end{equation}
Thus the strict \(\kappa\)-column form applies when
\(n\leq\kappa(p-1)\).  For a longer query, the number of initial
partition columns must be increased, or separated from the number of
soundness-amplification lanes; using only \(\kappa\) longer columns
would invalidate the characteristic-overflow argument below.
Pad the query with a fixed table value and reshape it columnwise as
\begin{equation}
    \bm V
      :=\concat{\bm v^{(1)},\ldots,\bm v^{(\kappa)}}
      \in\Field^{m\times\kappa}.
\label{eq:packed-logup-query-matrix}
\end{equation}
Padding with a valid table value does not change whether the original
query satisfies the lookup.

Column \(j\) has its own multiplicity vector.  The honest value at table
entry \(t_r\) is
\begin{equation}
    M_{r,j}
      :=\#\{i\in[m]:V_{i,j}=t_r\},
    \qquad
    \bm M\in\Field^{T\times\kappa}.
\label{eq:packed-logup-column-multiplicities}
\end{equation}
Because every column has length less than \(p\), each integer in
\eqref{eq:packed-logup-column-multiplicities} embeds into \(\Field\)
without wraparound.  The prover commits once to \(\bm M\); it need not
separately prove the counting formula.  Instead, column \(j\) is bound
to the rational identity
\begin{equation}
    F_j(X)
      :=
      \sum_{i=1}^{m}\frac{1}{X+V_{i,j}}
      -
      \sum_{r=1}^{T}\frac{M_{r,j}}{X+t_r}
      \stackrel{?}{=}0.
\label{eq:packed-logup-column-identity}
\end{equation}
If the column is valid, its honest multiplicities make this an identity
of rational functions.  Conversely, an invalid value produces a pole
that does not occur on the table side.

\paragraph{Combining the column identities.}
Assign the \(\kappa\) columns distinct addresses
\[
    \bm s^{(1)},\ldots,\bm s^{(\kappa)}
      \in\{0,1\}^{h}.
\]
For \(\bm y\in\Field^h\), define
\begin{equation}
    \bm\eta_{\bm y}
      :=
      \left(
        \bm y^{\bm s^{(1)}},\ldots,
        \bm y^{\bm s^{(\kappa)}}
      \right)^\top,
    \qquad
    \bm y^{\bm s}:=\prod_{q=1}^{h}y_q^{s_q}.
\label{eq:packed-logup-column-weights}
\end{equation}
The packed rational error is
\begin{equation}
    F_{\bm y}(X)
      :=\sum_{j=1}^{\kappa}
        (\bm\eta_{\bm y})_j F_j(X).
\label{eq:packed-logup-packed-rational-error}
\end{equation}
The monomials associated with the occupied addresses are distinct.
Hence, if at least one \(F_j\) is nonzero, then
\(F_{\bm y}(X)\) is a nonzero polynomial in \(\bm y\) whose
coefficients are rational functions of \(X\).

\paragraph{Independent packed evaluations.}
Let
\[
    \mathbb E:=\Field[\vartheta]/(\vartheta^2-\nu)
      \cong\mathbb F_{p^2},
\]
where \(\nu\in\Field\) is a quadratic non-residue.  After \(\bm V\)
and \(\bm M\) are bound, the verifier samples \(\kappa\) independent
pairs
\begin{equation}
    \left(x^{(c)},\bm y^{(c)}\right)
      \ChooseRandom
      (\mathbb E\setminus\Field)\times\Field^h,
    \qquad c\in[\kappa],
\label{eq:packed-logup-initial-challenges}
\end{equation}
and writes \(\bm\eta^{(c)}:=\bm\eta_{\bm y^{(c)}}\).  The intended
child claim in lane \(c\) is
\[
    F_{\bm y^{(c)}}(x^{(c)})=0.
\]

\paragraph{Creating the \(\kappa\) child columns.}
For every row \(i\in[m]\) and lane \(c\in[\kappa]\), define
\begin{equation}
    a_i^{(c)}
      :=
      \sum_{j=1}^{\kappa}
        \frac{\eta_j^{(c)}}{x^{(c)}+V_{i,j}},
    \qquad
    \bm A
      :=\concat{\bm a^{(1)},\ldots,\bm a^{(\kappa)}}
      \in\mathbb E^{m\times\kappa}.
\label{eq:packed-logup-child-columns}
\end{equation}
Although every child uses a different \(x^{(c)}\), the matrix
\(\bm A\) has only \(m\kappa=n+O(\kappa)\) extension-field entries.
In particular, it is not a \(\kappa\)-fold copy of the query.

The entries of \(\bm A\) are linked directly to the already-bound
matrix \(\bm V\).  Since \(x^{(c)}\notin\Field\), every denominator is
nonzero, and \eqref{eq:packed-logup-child-columns} is equivalent to the
polynomial relation
\begin{equation}
\boxed{
    a_i^{(c)}
      \prod_{j=1}^{\kappa}(x^{(c)}+V_{i,j})
    =
      \sum_{j=1}^{\kappa}
        \eta_j^{(c)}
        \prod_{q\in[\kappa]\setminus\{j\}}
          (x^{(c)}+V_{i,q}).}
\label{eq:packed-logup-packed-reciprocal}
\end{equation}
Thus the prover commits only to the combined values
\(a_i^{(c)}\), rather than to all
\(\kappa^2m\) individual inverses.  Relation
\eqref{eq:packed-logup-packed-reciprocal} has total degree
\(\kappa+1\) over \(\mathbb E\); equivalently, it gives two base-field
relations after expanding in the basis \((1,\vartheta)\).

The table side is reduced before any table-sized child vector is
formed.  Define
\begin{equation}
    \beta^{(c)}
      :=
      \sum_{r=1}^{T}
        \frac{(\bm M\bm\eta^{(c)})_r}
             {x^{(c)}+t_r}
      =
      \bm q_{x^{(c)}}^\top
      \bm M\bm\eta^{(c)},
    \qquad
    (\bm q_x)_r:=\frac1{x+t_r}.
\label{eq:packed-logup-table-opening}
\end{equation}
The vector \(\bm q_x\) is public.  Hence each \(\beta^{(c)}\) is a
rank-one linear opening of the single committed matrix \(\bm M\), not
a commitment to a \(T\times\kappa\) inverse matrix.  Its two base-field
coordinates are proved by the packed inner-product primitive.

By construction,
\begin{equation}
    F_{\bm y^{(c)}}(x^{(c)})
      =\bm1^\top\bm a^{(c)}-\beta^{(c)}.
\label{eq:packed-logup-child-sum}
\end{equation}
Consequently, the \(\kappa\) rational identities have become the
\(\kappa\) linear claims
\begin{equation}
    \boxed{
    \bm\beta
      \EqOrNot
      \bm A^\top\bm1,}
    \qquad
    \bm\beta
      :=(\beta^{(1)},\ldots,\beta^{(\kappa)})^\top.
\label{eq:packed-logup-initial-child-state}
\end{equation}

\paragraph{Recursive \(\kappa\)-to-\(\kappa\) repacking.}
Pad \(\bm A\) with zero rows to a power-of-two length.  At any later
round, the state is a matrix
\(\bm A\in\mathbb E^{m'\times\kappa}\) and a vector
\(\bm c\in\mathbb E^\kappa\) with the claimed invariant
\begin{equation}
    \bm c\EqOrNot\bm A^\top\bm1.
\label{eq:packed-logup-linear-invariant}
\end{equation}
Initially, \(\bm c=\bm\beta\).  Split the rows as
\[
    \bm A=
      \begin{pmatrix}\bm A_U\\\bm A_D\end{pmatrix}.
\]
The prover sends vectors \(\widehat{\bm u},\widehat{\bm d}
\in\mathbb E^\kappa\), purportedly equal to
\(\bm A_U^\top\bm1\) and \(\bm A_D^\top\bm1\), and the verifier checks
\begin{equation}
    \widehat{\bm u}+\widehat{\bm d}\EqOrNot\bm c.
\label{eq:packed-logup-parent-sum-check}
\end{equation}

After this message is bound, the verifier samples \(\kappa\) fresh
pairs
\[
    (\bm y_1,z_1),\ldots,(\bm y_\kappa,z_\kappa)
      \ChooseRandom\Field^h\times\Field.
\]
Set
\[
    \bm Y
      :=\begin{pmatrix}
          \bm\eta_{\bm y_1}&\cdots&\bm\eta_{\bm y_\kappa}
        \end{pmatrix},
    \qquad
    \bm Z:=\diag(z_1,\ldots,z_\kappa).
\]
The next \(\kappa\) columns and claims are
\begin{equation}
\boxed{
\begin{aligned}
    \bm A'
      &:=\bm A_U\bm Y+\bm A_D\bm Y\bm Z,\\
    \bm c'
      &:=\bm Y^\top\widehat{\bm u}
        +\bm Z\bm Y^\top\widehat{\bm d}.
\end{aligned}}
\label{eq:packed-logup-recursive-repacking}
\end{equation}
Thus every round maps \(\kappa\) columns of length \(m'\) to
\(\kappa\) columns of length \(m'/2\).  When one row remains, the
verifier checks \eqref{eq:packed-logup-linear-invariant} by opening the
terminal row.  The initial \(\kappa\) pairs
\((x^{(c)},\bm y^{(c)})\) create the first \(\kappa\) columns, and every
subsequent collection of \(\kappa\) pairs
\((\bm y_c,z_c)\) creates the next \(\kappa\) columns.

\paragraph{Folding invariant and round soundness.}
For an honest round, let
\[
    \bm u^\star:=\bm A_U^\top\bm1,
    \qquad
    \bm d^\star:=\bm A_D^\top\bm1.
\]
Then \eqref{eq:packed-logup-recursive-repacking} satisfies
\[
    (\bm A')^\top\bm1
      =\bm Y^\top\bm u^\star
       +\bm Z\bm Y^\top\bm d^\star,
\]
so the invariant is preserved.

For soundness, write
\[
    \bm e_U:=\widehat{\bm u}-\bm u^\star,
    \qquad
    \bm e_D:=\widehat{\bm d}-\bm d^\star.
\]
If a parent claim is false while
\eqref{eq:packed-logup-parent-sum-check} accepts, then
\(\bm e_U+\bm e_D\ne\bm0\).  The error in a child sampled at
\((\bm y,z)\) is
\begin{equation}
    \Delta(\bm y,z)
      :=\bm\eta_{\bm y}^\top(\bm e_U+z\bm e_D).
\label{eq:packed-logup-fold-error}
\end{equation}
This is a nonzero polynomial of total degree at most \(h+1\): setting
\(z=1\) leaves the nonzero channel polynomial
\(\bm\eta_{\bm y}^\top(\bm e_U+\bm e_D)\).  Therefore all
\(\kappa\) independently sampled children become correct with
probability at most
\begin{equation}
    \epsilon_{\mathsf{linear\mbox{-}round}}
      \leq
      \left(\frac{h+1}{p}\right)^\kappa.
\label{eq:packed-logup-linear-round-soundness}
\end{equation}

\paragraph{Lookup soundness.}
Suppose that
\(v^\star\notin\{t_1,\ldots,t_T\}\) occurs in the query, and let
\(n_j^\star\) be the number of its occurrences in column \(j\).  By
\eqref{eq:packed-logup-short-column-condition},
\[
    0\leq n_j^\star<p,
    \qquad
    (n_1^\star,\ldots,n_\kappa^\star)\ne\bm0.
\]
Its residue in the packed rational function is
\begin{equation}
    R_{v^\star}(\bm y)
      :=\sum_{j=1}^{\kappa}
        n_j^\star(\bm\eta_{\bm y})_j.
\label{eq:packed-logup-invalid-residue}
\end{equation}
This is a nonzero polynomial of total degree at most \(h\), so it
vanishes at a random \(\bm y\in\Field^h\) with probability at most
\(h/p\).  Conditioned on a nonzero residue, the resulting rational
function has a pole at \(-v^\star\) that cannot occur on the table side.
After its at most \(p\) distinct denominators are combined, its
numerator has degree at most \(p-1\).  A uniform
\(x\in\mathbb E\setminus\Field\) is therefore a root with probability
at most \(1/p\).  One initial pair fails to expose the invalid lookup
with probability at most
\[
    \delta
      :=\frac hp+\left(1-\frac hp\right)\frac1p
      \leq\frac{h+1}{p}.
\]
Because the \(\kappa\) pairs in
\eqref{eq:packed-logup-initial-challenges} are independent, all initial
child claims are correct with probability at most \(\delta^\kappa\).
If \(L=\lceil\log_2m\rceil\), a union bound over the initial packing and
the subsequent linear folds gives
\begin{equation}
\boxed{
    \epsilon_{\mathsf{PackedLogUp}}
      \leq
      (L+1)
      \left(\frac{h+1}{p}\right)^\kappa
      +\epsilon_{\mathsf{link}}
      +\epsilon_{\mathsf{com}}.}
\label{eq:packed-logup-soundness}
\end{equation}
Here \(\epsilon_{\mathsf{link}}\) contains the soundness errors of the
packed reciprocal relation and the rank-one openings of \(\bm M\), and
\(\epsilon_{\mathsf{com}}\) contains commitment binding and terminal
opening errors.

\paragraph{Commitments, cost, and memory.}
The transcript order is
\begin{equation}
\begin{aligned}
  \mathsf{Com}(\bm v,\bm t)
    &\longrightarrow \mathsf{Com}(\bm M)
    \longrightarrow
      \{(x^{(c)},\bm y^{(c)})\}_{c=1}^{\kappa}\\
    &\longrightarrow \mathsf{Com}(\bm A)
    \longrightarrow
      \text{linking proofs and recursive folds}.
\end{aligned}
\label{eq:packed-logup-transcript-order}
\end{equation}
The table commitment is omitted when the table is public.  No
query-by-challenge or table-by-challenge inverse matrix is committed.
There are only two logical lookup-specific witness commitments: the
multiplicity matrix \(\bm M\), bound before the initial challenges, and
the combined child matrix \(\bm A\), bound afterward.  Since
\(\bm A\) is extension-valued, its logical commitment consists of two
base-field coordinate commitments.

An honest multiplicity matrix has at most \(m\kappa\) nonzero entries,
even though its ambient dimension is \(T\times\kappa\).  It can
therefore be kept in sparse column form and fed to a single commitment;
streaming the encoder changes the prover schedule but not the
commitment or proof size.  The matrix \(\bm A\) itself has
\(m\kappa=n+O(\kappa)\) extension-field entries.  Hence the peak
working memory, excluding the public table, is
\begin{equation}
    \boxed{O(n+\kappa^2)}
\label{eq:packed-logup-working-memory}
\end{equation}
base-field elements, up to the constant factor for representing
\(\mathbb E\).  If the commitment implementation requires a resident
dense encoding of \(\bm M\), its \(T\kappa\) storage must instead be
added to this bound; the algebraic packing does not by itself remove
that implementation cost.

The memory bound also assumes that
\eqref{eq:packed-logup-packed-reciprocal} is handled as a direct
degree-\((\kappa+1)\) packed gate, or by a product check whose
intermediate values are streamed.  Expanding every product into
quadratic gates and materializing all partial products would add
\(\Theta(n\kappa)\) witness entries and would forfeit the intended
space saving.  Thus this construction trades the large inverse trace
for one moderate-degree linking relation; its efficient realization is
part of \(\epsilon_{\mathsf{link}}\) and the corresponding linking cost.

Computing all entries in \eqref{eq:packed-logup-child-columns}, all
table openings in \eqref{eq:packed-logup-table-opening}, and all dense
column folds costs \(O(n\kappa)\) extension-field operations when
\(\bm M\) is stored sparsely.  In each recursive round,
\(\widehat{\bm d}=\bm c-\widehat{\bm u}\) may be omitted, so the prover
sends \(\kappa\) extension-field elements.  The \(L\) linear folds
therefore contribute \(2\kappa L\) base-field elements, in addition to
the packed reciprocal proof, the rank-one openings of \(\bm M\), and
the terminal commitment opening.  If the two base-field coordinates of
the table openings are instantiated with Multiplicative Packed
LiteBullet, they add at most
\[
    2(3\kappa^2-\kappa)
      \left\lceil\log_2(T\kappa)\right\rceil
\]
base-field elements before terminal openings.

Packed LogUp proves membership with unrestricted reuse of table
entries.  Capacitated lookup or exact multiset equality additionally
requires a range-safe constraint on the multiplicities.


\subsection{Packed Permutation}
\label{sec:packed-permutation}

\paragraph{State.}
Let
\[
    \bm A=
      \left(\bm a^{(1)}\,\middle|\,\cdots\,\middle|\,
      \bm a^{(\kappa)}\right),
    \qquad
    \bm B=
      \left(\bm b^{(1)}\,\middle|\,\cdots\,\middle|\,
      \bm b^{(\kappa)}\right)
      \in\Field^{m\times\kappa},
\]
and write \(N:=m\kappa\) for the total number of entries.  We prove
the \(\kappa\) independent statements
\begin{equation}
    \bm b^{(c)}
      \text{ is a permutation of }
    \bm a^{(c)},
    \qquad c\in[\kappa].
\label{eq:packed-permutation-column-statements}
\end{equation}
The permutations may differ across columns, and repeated values are
allowed.  We assume throughout that
\begin{equation}
    m<p.
\label{eq:packed-permutation-channel-bound}
\end{equation}
This condition is needed only to interpret field-valued multiplicities
as integer multiplicities: every difference lies in \([-m,m]\), so a
nonzero difference cannot vanish modulo \(p\).

\paragraph{Power-sum representation.}
Set \(\bm X_A:=\bm A\) and \(\bm X_B:=\bm B\).  For
\(\sigma\in\{A,B\}\) and \(d=0,\ldots,p-1\), define the virtual power
tables
\begin{equation}
    \bm P_{\sigma,d}:=\bm X_\sigma^{\circ d}
      \in\Field^{m\times\kappa},
    \qquad
    \bm P_{\sigma,0}:=\bm 1_{m\times\kappa},
\label{eq:packed-permutation-power-tables}
\end{equation}
Stack their column sums by degree:
\begin{equation}
    \bm M_\sigma
      :=
      \begin{pmatrix}
        \bm 1_m^\top\bm P_{\sigma,0}\\
        \bm 1_m^\top\bm P_{\sigma,1}\\
        \vdots\\
        \bm 1_m^\top\bm P_{\sigma,p-1}
      \end{pmatrix}
      \in\Field^{p\times\kappa}.
\label{eq:packed-permutation-moment-matrix}
\end{equation}

For column \(c\), let \(\bm h^{(c)}\in\Field^p\) be its multiplicity
residue, indexed by \(x\in\Field\):
\[
    h_x^{(c)}
      :=
      \#\{i:A_{i,c}=x\}
      -\#\{i:B_{i,c}=x\}.
\]
Let \(\bm V=(x^d)_{x\in\Field,\,0\leq d<p}\) be the \(p\)-by-\(p\)
Vandermonde matrix over all elements of \(\Field\).  The \(c\)-th
column of the moment residue satisfies
\begin{equation}
    (\bm M_A-\bm M_B)[:,c]
      =\bm V^\top\bm h^{(c)}.
\label{eq:packed-permutation-vandermonde-link}
\end{equation}
The evaluation points indexing the rows of \(\bm V\) are distinct, so
\(\bm V\) is invertible.  Together with \(m<p\), this gives the exact
equivalence
\begin{equation}
    \boxed{
    \bm M_A=\bm M_B
      \quad\Longleftrightarrow\quad
    \bm b^{(c)}\text{ is a permutation of }\bm a^{(c)}
      \text{ for every }c\in[\kappa].
    }
\label{eq:packed-permutation-moment-equivalence}
\end{equation}

\paragraph{Random-polynomial fingerprint.}
Let \(s\) denote the number of soundness lanes.  After the commitments
to \(\bm A\) and \(\bm B\) are fixed, Fiat--Shamir generates
independent uniform coefficient matrices
\[
    \bm R^{(1)},\ldots,\bm R^{(s)}
      \ChooseRandom\Field^{p\times\kappa}.
\]
Lane \(t\) checks the scalar fingerprint
\begin{equation}
    \boxed{
    0
      \EqOrNot
      \Phi_t
      :=
      \left\langle
        \bm R^{(t)},\bm M_A-\bm M_B
      \right\rangle.
    }
\label{eq:packed-permutation-random-fingerprint}
\end{equation}
Equivalently, the coefficients in column \(c\) define the random
polynomial
\begin{equation}
    f_{t,c}(X)
      :=
      \sum_{d=0}^{p-1}R^{(t)}_{d,c}X^d,
\label{eq:packed-permutation-random-polynomial}
\end{equation}
and
\begin{equation}
    \Phi_t
      =
      \sum_{c=1}^{\kappa}\sum_{i=1}^{m}
      \left(
        f_{t,c}(A_{i,c})-f_{t,c}(B_{i,c})
      \right).
\label{eq:packed-permutation-polynomial-evaluation}
\end{equation}
Using a full random coefficient matrix, rather than one coefficient
vector shared by all columns, is essential: every soundness lane tests
the complete column-residue matrix through an independent projection.

\paragraph{Reducing the virtual power tables.}
The prover need not commit \(\bm M_A\), \(\bm M_B\), or the power
tables.  Their entries are generated by the degree-two recurrence
\begin{equation}
\begin{aligned}
    \bm P_{A,d+1}&=\bm P_{A,d}\circ\bm A,\\
    \bm P_{B,d+1}&=\bm P_{B,d}\circ\bm B,
\end{aligned}
\qquad d=1,\ldots,p-2,
\label{eq:packed-permutation-power-recurrence}
\end{equation}
with public layer \(\bm P_{\sigma,0}=\bm 1\) and committed layer
\(\bm P_{A,1}=\bm A\), \(\bm P_{B,1}=\bm B\).

Write \((\bm r_d^{(t)})^\top:=\bm R^{(t)}[d,:]\) and define the public
weight matrix
\[
    \bm W_d^{(t)}:=\bm 1_m(\bm r_d^{(t)})^\top
      \in\Field^{m\times\kappa}.
\]
Then the initial claim in lane \(t\) is
\begin{equation}
    \Phi_t
      =
      \sum_{d=0}^{p-1}
      \left\langle
        \bm W_d^{(t)},
        \bm P_{A,d}-\bm P_{B,d}
      \right\rangle.
\label{eq:packed-permutation-weighted-powers}
\end{equation}
For either operand and any public weight matrix \(\bm W\), one power
step has the inner-product form
\begin{equation}
    \left\langle\bm W,\bm P_{\sigma,d+1}\right\rangle
      =
      \left\langle
        \operatorname{vec}(\bm W\circ\bm P_{\sigma,d}),
        \operatorname{vec}(\bm X_\sigma)
      \right\rangle.
\label{eq:packed-permutation-power-inner-product}
\end{equation}
Stacking \eqref{eq:packed-permutation-power-inner-product} over
\(t\in[s]\) gives one \(s\)-channel Packed LiteBullet instance for
that degree.  Its returned linear claims on \(\bm P_{\sigma,d}\) are
re-packed with the degree-\(d\) contribution in
\eqref{eq:packed-permutation-weighted-powers} before the next step.
Iterating from \(d=p-1\) to \(d=1\) leaves only public claims on
\(\bm P_{\sigma,0}\) and dense linear claims on the committed matrices
\(\bm A\) and \(\bm B\).  The latter are accumulated and discharged by
the dense-linear-to-MLE link from Section~\ref{sec:packed-litebullet}.

\paragraph{Commitments and transcript order.}
The only persistent witness commitments are
\begin{equation}
    \boxed{
      \mathsf{Com}_{\Field}(\bm A),
      \qquad
      \mathsf{Com}_{\Field}(\bm B).
    }
\label{eq:packed-permutation-commitments}
\end{equation}
The coefficient matrices are derived from the Fiat--Shamir transcript
after these commitments and can be generated from their transcript
seeds on demand.  No bit decomposition, permutation matrix, moment
matrix, or power-table commitment is introduced.  Fresh re-packing
challenges are sampled only after the preceding Packed LiteBullet
messages have been absorbed.

\paragraph{Soundness.}
Suppose that at least one column statement is false.  By
\eqref{eq:packed-permutation-moment-equivalence}, the fixed residue
\[
    \bm E:=\bm M_A-\bm M_B
\]
is nonzero.  For a uniform \(\bm R^{(t)}\), the inner product
\(\langle\bm R^{(t)},\bm E\rangle\) is uniform in \(\Field\), and hence
\[
    \Pr[\Phi_t=0\mid\bm E\neq\bm0]=\frac1p.
\]
The \(s\) coefficient matrices are independent, so
\begin{equation}
    \boxed{
    \epsilon_{\mathsf{perm\mbox{-}fingerprint}}
      =p^{-s}.
    }
\label{eq:packed-permutation-soundness}
\end{equation}
Taking \(s=\kappa\) gives \(p^{-\kappa}\).  The complete error bound is
\[
    \epsilon_{\mathsf{PackedPerm}}
      \leq
      p^{-s}
      +\epsilon_{\mathsf{power}}
      +\epsilon_{\mathsf{PackedLiteBullet}}
      +\epsilon_{\mathsf{link}}
      +\epsilon_{\mathsf{com}},
\]
where \(\epsilon_{\mathsf{power}}\) collects the fresh re-packing
errors across the \(p-2\) power steps.

The bound \(m<p\) is essential.  If a multiplicity residue is a
nonzero multiple of \(p\), then it is already zero in \(\Field\), so
neither random coefficients nor an extension of \(\Field\) can recover
the lost integer count.  Longer columns require range-safe count limbs
or a preliminary partition into channels of length smaller than \(p\).

\paragraph{Cost and memory.}
Constructing all virtual powers and their column sums once costs
\begin{equation}
    O(pN)
\label{eq:packed-permutation-trace-cost}
\end{equation}
base-field operations.  Forming \(s\) fingerprints from the moment
matrix costs another \(O(sp\kappa)\) operations.  Proving all
\(p-2\) degree transitions with \(s\)-channel Packed LiteBullet costs
\(O(spN)\) prover operations up to the logarithmic factors of the
inner-product reduction.  Thus, for \(s=\kappa\), the proof-generation
cost is \(O(p\kappa N)\) up to those factors.

The proof contains \(p-2\) packed power-step reductions, followed by
the terminal dense openings.  Hence its size is \(p-2\) times the size
of one \(s\)-channel Packed LiteBullet instance, up to the final
opening terms.  The \(sp\kappa\) random coefficients are generated
from Fiat--Shamir seeds and are not transmitted in the proof.

Materializing every power layer requires \(O(pN)\) field elements.
Instead, the prover may retain one \(m\)-by-\(\kappa\) power layer and
the \(p\)-by-\(\kappa\) moment matrix, for
\begin{equation}
    O(N+p\kappa)
\label{eq:packed-permutation-streaming-memory}
\end{equation}
working memory.  Layers needed later by the recursive proof are
recomputed from \(\bm A\) and \(\bm B\), giving the usual time--memory
tradeoff.  The random-polynomial construction therefore removes the
bit matrix, the grand product, and extension-field state, at the price
of \(p-2\) degree-two power transitions.  This tradeoff is attractive
precisely when \(p\) is small.

\subsection{Polynomial Commitments over Small Fields}
\label{sec:pcs-small-field}

We instantiate the commitment layer with rank-one QA codes over
\(\Field=\mathbb F_p\).

\begin{definition}[Walsh--Hadamard transform]
Let \(N\) be a power of two and assume that \(N\) is invertible in
\(\Field\).  Define the \(N\times N\) Hadamard matrix recursively by
\[
    \bm H_1:=(1),
    \qquad
    \bm H_N
    :=
    \begin{pmatrix}
        \bm H_{N/2} & \bm H_{N/2}\\
        \bm H_{N/2} & -\bm H_{N/2}
    \end{pmatrix}
    =
    \begin{pmatrix}
        1&1\\
        1&-1
    \end{pmatrix}^{\otimes\log_2N}.
\]
For \(\bm a\in\Field^N\), its Walsh--Hadamard transform is the matrix--vector
product
\[
    \mathsf{WHT}_N(\bm a):=\bm H_N\bm a.
\]
Since \(\bm H_N^2=N\bm I_N\), it satisfies
\[
    \mathsf{WHT}_N(\mathsf{WHT}_N(\bm a))=N\bm a,
    \qquad
    \mathsf{WHT}_N^{-1}=N^{-1}\mathsf{WHT}_N.
\]
We omit the subscript \(N\) when the vector length is clear.
\end{definition}

\begin{lemma}[MLE reduction through the Walsh--Hadamard transform]
\label{lemma:wht-mle}
Let \(\bm A\in\Field^{N\times M}\), and define the column-wise
Walsh--Hadamard transform by
\[
    \mathsf{WHT}_N(\bm A):=\bm H_N\bm A.
\]
For
\(
    \bm l=\tensorspan{\bm x_l}\in\Field^N
\)
and
\(
    \bm r=\tensorspan{\bm x_r}\in\Field^M
\),
we have
\[
    \bm l^\top\mathsf{WHT}_N(\bm A)\bm r
    =
    \widetilde{\bm l}^{\top}\bm A\bm r,
    \qquad
    \widetilde{\bm l}
    :=
    \bigotimes_{i=1}^{\log_2N}(1,1-2x_{l,i}).
\]
Moreover, if \(x_{l,i}\ne1\) for every \(i\), define
\[
    \lambda_l
    :=
    \prod_{i=1}^{\log_2N}(2-2x_{l,i}),
    \qquad
    \bm x_l'
    :=
    \frac{\bm{1}-2\bm x_l}{2\bm{1}-2\bm x_l}.
\]
Then
\[
    \widetilde{\bm l}
    =
    \lambda_l\tensorspan{\bm x_l'},
\]
and hence
\[
    \bm l^\top\mathsf{WHT}_N(\bm A)\bm r
    =
    \lambda_l\,
    \tensorspan{\bm x_l'}^\top\bm A\bm r.
\]
\end{lemma}

\begin{proof}
Since \(\bm H_N\) is symmetric,
\[
    \bm l^\top\bm H_N\bm A\bm r
    =
    (\bm H_N\bm l)^\top\bm A\bm r.
\]
For every \(i\),
\[
    \begin{pmatrix}1&1\\1&-1\end{pmatrix}
    \begin{pmatrix}1-x_{l,i}\\x_{l,i}\end{pmatrix}
    =
    \begin{pmatrix}1\\1-2x_{l,i}\end{pmatrix}.
\]
Taking the tensor product over all \(i\) gives the expression for
\(\widetilde{\bm l}\).  Finally,
\[
    (1,1-2x_{l,i})
    =
    (2-2x_{l,i})
    \left(
        1-\frac{1-2x_{l,i}}{2-2x_{l,i}},
        \frac{1-2x_{l,i}}{2-2x_{l,i}}
    \right),
\]
which proves the normalized form.
\end{proof}

\begin{definition}[Quasi-Abelian code {\cite[Theorem~4.3]{iacr/LiLXYY26}}]
Let \(c=r\) be the index of the QA code, so its rate is \(1/c=1/r\), and let
\(\bm a_1,\ldots,\bm a_c\in\Field^{k_j}\) be the coefficient vectors of
its defining group-algebra elements.  Define
\[
    \bm\alpha_i
    :=
    k_j^{-1}\mathsf{WHT}(\bm a_i),
    \qquad i\in[c].
\]
For a message \(\bm m\in\Field^{k_j}\), its encoding is
\[
    \Enc(\bm m)
    :=
    \concat{
        \mathsf{WHT}\!\left(
            \mathsf{WHT}(\bm m)\circ\bm\alpha_1
        \right),
        \ldots,
        \mathsf{WHT}\!\left(
            \mathsf{WHT}(\bm m)\circ\bm\alpha_c
        \right)
    },
\]
where \(\circ\) denotes component-wise multiplication.
In the systematic variant, one sets \(\bm a_1\) to the multiplicative
identity, in which case the first output block equals \(\bm m\).
\end{definition}

\begin{fact}[Relative distance of the QA code
{\cite[Corollary~3.21]{iacr/LiLXYY26}}]
\label{fact:qa-distance}
Let \(k_j=2^n\), let \(\Field=\mathbb F_p\) for an odd prime \(p\), and define
\[
    g_p(\delta)
    :=
    1-\delta\log_p(p-1)
      +\delta\log_p\delta
      +(1-\delta)\log_p(1-\delta).
\]
For a target relative distance \(\delta\), suppose that
\[
    \epsilon
    :=
    g_p(\delta)-\frac{1+\log_p k_j}{c}
    >0.
\]
Then a uniformly sampled rank-one QA code of index \(c\) over
\(\mathbb F_p[\mathbb Z_2^n]\) has rate \(1/c\), length \(ck_j\), and
relative distance at least \(\delta\) with probability at least
\[
    \max\left\{
        1-\frac{c(c-1)k_j}{2p^2}
          -\frac{
              p^{-\left\lceil\frac{c-1}{c\delta}\right\rceil c\epsilon}
          }{
              (1-p^{-c\epsilon})(p-1)
          },
        \;
        1-\frac{ck_j}{p}
          -\frac{
              p^{-\left\lceil\frac{1}{\delta}\right\rceil c\epsilon}
          }{
              (1-p^{-c\epsilon})(p-1)
          }
    \right\}.
\]
The same distance analysis applies to the systematic construction above
by~\cite[Theorem~4.4]{iacr/LiLXYY26}.
\end{fact}

\medskip
\noindent
\textbf{Instantiating \methodname.}
We call the resulting instantiation \qafold.
At every level \(j\), \qafold instantiates \(\Code_j\) with a rank-one QA code
over \(\Field[\mathbb Z_2^{\log_2 k_j}]\) and uses its WHT encoder as
\(\Enc_j\).  Its generator matrix is
\[
    \bm G_j
    =
    \begin{pmatrix}
        \bm H_{k_j}\diag(\bm\alpha_1)\bm H_{k_j}\\
        \vdots\\
        \bm H_{k_j}\diag(\bm\alpha_c)\bm H_{k_j}
    \end{pmatrix}
    \in
    \Field^{ck_j\times k_j}.
\]
The fixed matrix used to reduce MLE claims on \(\bm G_j\) is
\[
    \bm H_j
    :=
    \concat{\bm\alpha_1,\ldots,\bm\alpha_c}
    \in
    \Field^{k_j\times c}.
\]
Applying Lemma~\ref{lemma:wht-mle} to the two WHT layers reduces every
MLE claim on \(\bm G_j\) to an MLE claim on \(\bm H_j\), together with
the corresponding scalar normalization factors.  Thus, in the notation
of Section~\ref{sec:decider}, the setup commitment at level \(j\)
authenticates \(\concat{\bm T_j,\bm H_j}\), while \(\bm S_j\) is the empty
matrix and no dynamic commitment is needed.  The carry commitment to
\(\bm W_j\) and the recursive code switch between
\(\widehat{\Code}_j\) and \(\widehat{\Code}_{j-1}^{t}\) are unchanged.
Consequently, the same recursive opening procedure applies until the
remaining MLE claims on \(\bm G_0\) are evaluated directly by the verifier.


\subsection{Packing Independent Circuits}
\label{sec:independent-circuits}
Packed LiteBullet requires the \(\kappa\) claims to share only the base
field and the current vector length.  Their circuits may otherwise have
unrelated matrices, wiring, constants, public inputs, and witnesses.
Circuits of different sizes are padded to a common length before
packing.

\medskip \noindent
{\bf Allocating Circuits to the Channels.}
Each computation may occupy one or more of the \(\kappa\) channels.
A single channel suffices for one inner-product claim, whereas a
computation containing several inner-product checks may use several
channels.  The \(\kappa\) channels count only occupied positions, so
the address space of size \(2^h\) introduces no dummy claims.  When
the channels are used only for soundness amplification, they instead
contain copies of the same initial claim.

The columns of \(\bm A\) and \(\bm B\) are formed independently from
these claims.  In the quadratic form
\(\bm\eta_{\bm y}^\top\bm C^\star\bm\eta_{\bm y}\), column pair
\((r,s)\) is assigned the monomial
\(\bm y^{\bm s^{(r)}+\bm s^{(s)}}\).  The diagonal pair \((i,i)\)
occupies the distinct corner \(\bm y^{2\bm s^{(i)}}\), while no
cross-column pair can occupy that corner.
Equation~\eqref{eq:wht-interpolation} therefore isolates each parent
claim independently.  The reciprocal \(z\)-weights place the two
matching \(U/D\)-half products in the \(z^0\) coefficient and the
cross-half products in the \(z\) and \(z^{-1}\) coefficients.
Finally, \(\bm Y\) and \(\bm Z\) fold the witness columns and their
\(U/D\) halves at the same challenge points used to define the child
values in \eqref{eq:wht-state-update}.  The folding invariant then
ensures that every honest child claim remains the inner product of its
correspondingly folded witness columns.
Consequently, batching does not require the circuits to share their
arithmetic structure; they share only the field, dimensions, and
folding schedule.

\input{preliminary}

\clearpage

\section*{Acknowledgement}
The authors acknowledge the use of a generative AI tool for grammar, spelling, and minor language editing during manuscript preparation. No scientific content in this paper was generated by AI.

\bibliographystyle{splncs04}
\bibliography{ref}

\appendix
\section{Security Definitions}
\label{appx:def}

\subsection{Completeness and Soundness}
\label{appx:soundness}
\begin{definition}[Proof/Argument/Argument of Knowledge]
\label{def:proofArgumentAoK}
A protocol consists of algorithms $(\Setup,\Prover,\Verifier)$.
We distinguish proofs, arguments, and arguments of knowledge by
requiring soundness, computational soundness, and knowledge
soundness, respectively. Throughout, we require perfect completeness.
\begin{itemize}
	\item $\Setup$:
	On input $1^\lambda$, output public parameters $\pp$.
	\item Prover $\Prover$:
	On input $\pp$ and a statement $u$, produce a transcript $\pi$
    for membership in $\Language$. For a relation $\Relation$,
    the prover additionally receives a witness $w$.
	\item Verifier $\Verifier$:
	On input $\pp$, a statement $u$, and a transcript $\pi$,
    output {\sf true} or {\sf false}.
\end{itemize}
Let $\PVinteract$ denote the interaction between $\Prover$ and $\Verifier$.
\end{definition}

\begin{definition}[Perfect Completeness]
	\label{def:perfect-completeness}
The protocol $(\Setup,\Prover,\Verifier)$ for a relation \(\Relation\)
has perfect completeness if, for every \((u,w)\in\Relation\):
\[
    \Pr\left[
        \Verifier(\pp,u,\pi)=\TRUE
        \;\middle|\;
        \pp\leftarrow\Setup(1^\lambda),\;
        \pi\leftarrow\Prover(\pp,u,w)
    \right]
    =1,
\]
where the probability is over the randomness of
\(\Setup\), \(\Prover\), and \(\Verifier\).
\end{definition} 

\begin{definition}[Soundness~\cite{iacr/CanettiCHLRR18}]
	\label{def:soundness}
	A protocol $(\Setup,\Prover,\Verifier)$ for a language
    $\Language$ or relation $\mathcal R$ has soundness error
    $\kappa^\mathsf{s}(u)\colon\mathbb N\mapsto[0,1]$ if every
    computationally unbounded interactive prover $\Prover^*$ satisfies:
	\begin{equation*}
	% \label{eq:soundnessExpectedTime}
	\Pr
		\left(
			\left.
			\begin{aligned}
				& \Verifier(\pp, u, \pi) = \TRUE \\
			\end{aligned}
			\right.
			\left|
			\begin{aligned}
				& u \notin \Language~/~ \Relation;~ \pi \leftarrow (\Prover^* \interact \Verifier) (\pp, u) \\
			\end{aligned}
			\right.
		\right)
	\le \kappa^\mathsf{s}(u).
	\end{equation*}
    The protocol is sound if $\kappa^\mathsf{s}(u)$ is negligible.
\end{definition}

Restricting the adversarial prover to probabilistic polynomial time
gives computational soundness.
\begin{definition}[Computational Soundness]
	\label{def:ComputationalSoundness}
	A protocol $(\Setup,\Prover,\Verifier)$ for a language
    $\Language$ or relation $\mathcal R$ has computational soundness
    error $\kappa^\mathsf{cs}(u)\colon\mathbb N\mapsto[0,1]$ if
    every PPT prover $\Prover^*$ satisfies:
	\begin{equation*}
	% \label{eq:soundnessExpectedTime}
	\Pr
		\left(
			\left.
			\begin{aligned}
				& \Verifier(\pp, u, \pi) = \TRUE \\
			\end{aligned}
			\right.
			\left|
			\begin{aligned}
				& u \notin \Language~/~\Relation;~ \pi \leftarrow (\Prover^* \interact \Verifier)(\pp, u) \\
			\end{aligned}
			\right.
		\right)
	\le \kappa^\mathsf{cs}(u).
	\end{equation*}
    The protocol is computationally sound if
    $\kappa^\mathsf{cs}(u)$ is negligible.
\end{definition}
By definition, $\kappa^\mathsf{cs}(u)\leq\kappa^\mathsf{s}(u)$.

An argument of knowledge additionally requires extraction of a
witness~\cite{crypto/AttemaCK21}.
\begin{definition}[Knowledge Soundness]
	\label{def:knowledgeSoundness}
	A protocol $(\Setup,\Prover,\Verifier)$ for a relation $\Relation$
    has knowledge-soundness error
    $\kappa^\mathsf{ks}(u)\colon\mathbb N\mapsto[0,1]$ if there
    exists a knowledge extractor $\Extractor$ with the following
    guarantee. Given $u$ and black-box access to a prover $\Prover^*$,
    the extractor runs in expected time polynomial in $|u|$, counting
    each oracle query as one step, and satisfies:
	\begin{equation*}
	% \label{eq:soundnessExpectedTime}
	\Pr((u, \Extractor(\Prover^*, u) ) \in \Relation )
	\ge
	\epsilon(\Prover^*, u) - \kappa^\mathsf{ks}(u),
	\end{equation*}
	where $\epsilon(\Prover^*,u)$ is the acceptance probability of
    $\Prover^*$. The protocol is knowledge sound if
    $\kappa^\mathsf{ks}(u)$ is negligible.
\end{definition}

If a protocol is knowledge-sound with knowledge-soundness error
$\kappa^\mathsf{ks}(u)$, then it is also computationally sound with
error $\kappa^\mathsf{cs}(u) \le \kappa^\mathsf{ks}(u)$.
Otherwise, some false statement $u$ and PPT prover $\Prover^*$
would satisfy:
\[
	\Pr
	\left(
		\left.
		\begin{aligned}
			& \Verifier(\pp, u, \pi) = \TRUE \\
		\end{aligned}
		\right.
		\left|
		\begin{aligned}
			& u \notin \mathcal{R};~ \pi \leftarrow (\Prover^* \interact \Verifier)(\pp, u) \\
		\end{aligned}
		\right.
	\right)
	> \kappa^\mathsf{ks}(u).
\]
The extraction guarantee would then yield a witness $w$ with
$(u,w)\in\Relation$ with positive probability, contradicting the
falsity of $u$.


\subsection{Round-by-Round Soundness}
\label{appx:rbyDef}

\begin{definition}[Round-by-Round Soundness~\cite{iacr/CanettiCHLRR18}]
	\label{def:roundByRoundSoundness}
	Let $\Pi=(\Setup,\Prover,\Verifier)$ be a $(2\nu+1)$-message
    public-coin interactive proof for a relation $\Relation$.
    For a statement $u\in\{0,1\}^*$, write
    $\pi_j=(\alpha_0,x_1,\alpha_1,\ldots,\alpha_{j-1},x_j,\alpha_j)$
    for a transcript prefix containing the first $2j+1$ messages,
    where $j\in[0,\nu]$. Let $\mathsf H(u,\pi_j)$ denote the
    distribution of the verifier's next message.
    The protocol has round-by-round soundness error
    $\{\kappa_j^\mathsf{rbr}\}_{j\in[0,\nu]}$ if there exist
    deterministic, not necessarily efficient, functions
    $\{\mathsf{State}_j\}_{j\in[0,\nu]}$ mapping a statement and
    transcript prefix to $\mathsf{accept}$ or $\mathsf{reject}$,
    with the following properties:
		\begin{itemize}
			\item If $u \notin \Language$, then $\mathsf{State}_0(u, \emptyset) = \mathsf{reject}$, where $\emptyset$ denotes the empty transcript.
			\item If $\mathsf{State}_{j-1}(u,\pi_{j-1})=\mathsf{reject}$,
            then the next verifier challenge and prover response satisfy:
			\[
				\Pr\!\!
				\left(
					\mathsf{State}_j(u, \pi_j) = \mathsf{accept~}\!\!
					\left|~\!
						\begin{aligned}
							&\!\mathsf{State}_{j - 1}(u, \pi_{j - 1}) = \mathsf{reject}; \\
							& x_j \leftarrow \mathsf{H}(u, \pi_{j - 1}); \pi_j \leftarrow (\pi_{j - 1} \| x_j \| \alpha_j)\!
						\!\end{aligned}
					\right.
				\right)\!
				\!\le\!\kappa_j^\mathsf{rbr}.
			\]
			\item For any full transcript $\pi$ for which the verifier halts, if $\mathsf{State}_\nu(u, \pi) = \mathsf{reject}$ then $\Verifier(u, \pi) = \FALSE$.
		\end{itemize}
\end{definition}

The following lemma bounds the overall error by the sum of the
round-by-round errors.
\begin{lemma}[Round-by-Round Soundness Implies Soundness~\hspace{-1.58pt}\textup{\cite[Proposition~5.4]{iacr/CanettiCHLRR18}}]
\label{lemma:rbr}
	Let $\Pi$ be a $(2 \nu + 1)$-move interactive proof with round-by-round soundness error $\kappa_j^\mathsf{rbr}, j \in [\nu]$.
	Then $\Pi$ has soundness error at most
    $\sum_{j=1}^\nu\kappa_j^\mathsf{rbr}$.
\end{lemma}

\section{Deferred Proofs}
\label{appx:proof}
\label{appx:auxiliary-proofs}

This appendix proves the claims in Sections~\ref{sec:blocks}
and~\ref{sec:together} and the packed PCS claim in
Appendix~\ref{appx:pcs}, using the notation of each statement.

\subsection{Proof of Proposition~\ref{proposition:ip-complexity}}
\label{appx:proof-ip}

\begin{proof}
\noindent\textbf{Completeness.}
Expand the Laurent polynomial in Eq.~\eqref{eq:ip-laurent-polynomial}:
\[
    f(X)=
    \sum_{\tau=-(\kappa-1)}^{\kappa-1}f_\tau X^\tau,
    \qquad
    f_\tau=
    \sum_{\substack{j,j'\in[\kappa]\\j-j'=\tau}}
       \langle\bm a^{(j)},\bm b^{(j')}\rangle.
\]
Its constant coefficient is \(f_0=\langle\bm a,\bm b\rangle\).
Thus the auxiliary polynomials in Algorithm~\ref{alg:ip}
satisfy Eq.~\eqref{eq:ip-correction-identity} for a correct
inner-product claim. Evaluation at \(\ext x\neq0\) gives the
length-\(m\) claim in
Eq.~\eqref{eq:packed-inner-product-check}.

The terminal values of Fact~\ref{fact:ip} have the matrix form:
\[
\begin{aligned}
    \ext v_1
      &=\tensorspan{\ext{\bm y}}^\top\bm A\ext{\bm\rho}_A,
    &\ext{\bm\rho}_A
      &=(1,\ext x,\ldots,\ext x^{\kappa-1})^\top,\\
    \ext v_2
      &=\tensorspan{\ext{\bm y}}^\top\bm B\ext{\bm\rho}_B,
    &\ext{\bm\rho}_B
      &=(1,\ext x^{-1},\ldots,\ext x^{-(\kappa-1)})^\top.
\end{aligned}
\]
Both have tensor-structured row weights and fit the
packed-opening relation in Eq.~\eqref{eq:pcs-opening-relation}.

\smallskip\noindent\textbf{Soundness.}
Suppose \(c\neq\langle\bm a,\bm b\rangle\), and fix the
degree-\((\kappa-2)\) polynomials \(h_+,h_-\) before sampling
\(\ext x\). The Laurent residual
\(f(X)-c-Xh_+(X)-X^{-1}h_-(X^{-1})\) has a nonzero
constant coefficient. Multiplying by
\(X^{\kappa-1}\) gives a nonzero polynomial of degree at most
\(2\kappa-2\), so the evaluation check accepts with probability
at most \((2\kappa-2)/(|\ExtensionField|-1)\).
Conditioned on a false evaluated claim, Fact~\ref{fact:ip}
contributes at most \(2\log_2m/|\ExtensionField|\).
A union bound gives the stated error, excluding the terminal
packed openings.

\smallskip\noindent\textbf{Efficiency.}
Each ordered pair of distinct columns contributes a
length-\(m\) base-field inner product to the correction
coefficients, for \(m\kappa(\kappa-1)=N(\kappa-1)\) base-field
multiplications. Given the powers of \(\ext x\) and
\(\ext x^{-1}\), packing both matrices uses \(2N\,\be\).
Generating these powers and evaluating the corrections costs
\(O(\kappa)\,\ee\); Fact~\ref{fact:ip} adds
\((4m+O(\log m))\,\ee\).
Thus the total arithmetic cost is:
\[
    O(N\kappa)\,\bb+O(N)\,\be+O(m+\kappa)\,\ee
       =O(N\kappa+\kappa^3)\,\bb.
\]
Streaming the cross-column inner products retains only
\(O(\kappa)\) base-field correction coefficients.
The witness matrices and packed vectors use \(O(N)\)
base-field coordinates, and the challenge powers use
\(O(\kappa^2)\). An in-place inner-product reduction therefore
has peak memory \(O(N+\kappa^2)\Field\).
\end{proof}

\subsection{Proof of Proposition~\ref{proposition:hadamard-complexity}}
\label{appx:proof-hadamard}

\begin{proof}
\noindent\textbf{Completeness.}
The honest column message is
\(\ext{\bm c}^{\top}
=\ext{\bm l}^{\top}(\bm A\circ\bm B)\), so the scalar check
\(\ext c=\langle\ext{\bm c},\ext{\bm r}\rangle\) accepts.
At every interpolation node, the residual satisfies:
\[
    f(\gamma_j)
      =\langle\ext{\bm l},
          \bm a^{(j)}\circ\bm b^{(j)}\rangle-\ext c_j
      =0,
    \qquad j\in[\kappa].
\]
Hence \(f=\Psi h\) for a polynomial \(h\) of degree at most
\(\kappa-2\). Evaluating this identity at \(\ext x\) gives:
\[
    \langle\ext{\bm a},\ext{\bm b}\circ\ext{\bm l}\rangle
      =\ext c_x+\Psi(\ext x)h(\ext x).
\]
The terminal values of Fact~\ref{fact:ip} satisfy:
\[
\begin{aligned}
    \ext v_1
      &=\tensorspan{\ext{\bm y}}^\top
          \bm A\ext{\bm\rho}(\ext x),\\
    \ext v_2
      &=\bigl(\ext{\bm l}\circ\tensorspan{\ext{\bm y}}\bigr)^\top
          \bm B\ext{\bm\rho}(\ext x).
\end{aligned}
\]
Writing the supplied row weight as
\(\ext{\bm l}=\bigotimes_{j=1}^{\log_2m}
(\ext l_{j,0},\ext l_{j,1})^\top\), its product with the MLE
weight factors as:
\[
    \ext{\bm l}\circ\tensorspan{\ext{\bm y}}
      =\bigotimes_{j=1}^{\log_2m}
        \begin{pmatrix}
            \ext l_{j,0}(1-\ext y_j)\\
            \ext l_{j,1}\ext y_j
        \end{pmatrix}.
\]
Both terminal claims fit Eq.~\eqref{eq:pcs-opening-relation},
with Lagrange column weights and tensor-structured row weights.

\smallskip\noindent\textbf{Soundness.}
Suppose the weighted Hadamard claim is false but the scalar
check accepts. Then the column message differs from
\(\ext{\bm l}^{\top}(\bm A\circ\bm B)\) in at least one
coordinate, and \(f(\gamma_j)\neq0\) for some \(j\).
For every degree-\((\kappa-2)\) polynomial \(h\) fixed before
\(\ext x\), the residual \(f-\Psi h\) is therefore nonzero
and has degree at most \(2\kappa-2\).
A random \(\ext x\in\ExtensionField\setminus\mathcal H\)
makes this residual vanish with probability at most
\((2\kappa-2)/(|\ExtensionField|-\kappa)\).
Fact~\ref{fact:ip} contributes at most
\(2\log_2m/|\ExtensionField|\), giving:
\[
    \epsilon_{\mathsf{Had}}
      \leq
      \frac{2\kappa-2}{|\ExtensionField|-\kappa}
      +\frac{2\log_2m}{|\ExtensionField|}
      \leq
      \frac{2\kappa-2}{|\ExtensionField|-\kappa}
      +\frac{2\log_2N}{|\ExtensionField|}.
\]
The first inequality bounds Algorithm~\ref{alg:hadamard};
the second gives the stated bound. Terminal packed openings
are accounted for separately.
The argument applies unchanged when the interpolation nodes
lie in \(\mathbb K\subseteq\ExtensionField\).

\smallskip\noindent\textbf{Efficiency over the base field.}
First suppose \(\mathbb K=\Field\).
Computing \(\ext{\bm c}\) uses \(O(N)\,\bb+O(N)\,\be\).
At each of the \(\kappa-1\) additional interpolation nodes,
evaluating both witness encodings takes \(O(N)\,\bb\).
Pointwise multiplication and contraction with \(\ext{\bm l}\)
then take \(O(m)\,\bb+O(m)\,\be\).
Over all nodes, these operations cost
\(O(N\kappa)\,\bb+O(N)\,\be\).
Evaluating \(g_C\), interpolating the residual, and dividing by
\(\Psi\) contribute \(O(\kappa^2)\,\be\).

After \(\ext x\) is sampled, packing both matrices costs
\(O(N)\,\be\).
Expanding the row factors, forming
\(\ext{\bm b}\circ\ext{\bm l}\), and running Fact~\ref{fact:ip}
cost \(O(m)\,\ee\).
The terminal row factors require only componentwise
multiplication of their two-entry representations.
With precomputed Lagrange coefficients, the public query
and polynomial evaluations cost
\(O(\kappa^2)\,\be+O(\kappa)\,\ee\).
The total is therefore:
\[
    O(N\kappa)\,\bb
      +O(N+\kappa^2)\,\be
      +O(m+\kappa)\,\ee
      =O(N\kappa+\kappa^3)\,\bb.
\]

\smallskip\noindent\textbf{Intermediate-field tables.}
Let \([\mathbb K:\Field]=t>1\), with \(t\mid\kappa\) and
\(|\mathbb K|\geq2\kappa-1\).
Use the addition and multiplication tables specified in
Section~\ref{sec:hadamard}, charging each lookup as \(O(1)\).
Evaluating the witness encodings at all additional nodes
requires \(O(N\kappa)\) table operations.
In \(\mathbb K\)-coordinates, each extension-field element
has length \(\kappa/t\), so weighting and accumulating the
\(O(N)\) local products uses \(O(N\kappa/t)\) lookups.
The \(O(\kappa^2)\) mixed operations for \(g_C\), interpolation,
and division cost \(O(\kappa^3/t)\).
The remaining packing and inner-product steps retain their
\(O(N\kappa+\kappa^3)\,\bb\) bound. Thus the same online
time bound holds, excluding construction of the fixed tables.

\smallskip\noindent\textbf{Memory.}
Processing rows and sample points sequentially retains
\(O(\kappa)\) extension-field residual values, followed by
\(O(\kappa)\) extension-field coefficients after interpolation.
The packed vectors and inner-product reduction use \(O(m)\)
extension-field elements. Including the witness matrices,
this is \(O(N+\kappa^2)\) base-field coordinates.
The fixed interpolation data use \(O(\kappa^2t)\)
coordinates, and the tables use \(O(tq^{2t})\) when \(t>1\);
no field tables are needed when \(t=1\).
The peak memory is therefore
\(O\bigl(N+\kappa^2t+(t-1)q^{2t}\bigr)\Field\).
\end{proof}

\subsection{Proof of Proposition~\ref{proposition:permutation-complexity}}
\label{appx:proof-permutation}

\begin{proof}
\noindent\textbf{Completeness.}
The coefficient tree in
Eq.~\eqref{eq:packed-permutation-product-tree} computes the
products of \(2^\tau\) shifted witness entries in each column.
Evaluating its last layer at \(\ext z\) gives \(m\)
extension-field entries, whose product is
\(\mathsf{GrandProd}(\ext z\bm1+\bm a)\).
The two grand products agree whenever \(\bm a\) and \(\bm b\)
contain the same multiset of entries.

For a transition from level \(\tau+1\) to level \(\tau\),
Lagrange interpolation gives:
\[
    \bm g_{\tau,U}(Z,\gamma_j)
      \circ\bm g_{\tau,D}(Z,\gamma_j)
       =\bm C_{\tau+1}(Z)_{*,j},
    \qquad j\in[\kappa].
\]
For the honest column claims
\(\ext{\bm u}_{\tau+1}^\top
=\ext{\bm l}_{\tau+1}^\top\bm C_{\tau+1}(\ext z)\),
the residual \(f(X)\) therefore vanishes on \(\mathcal H\).
Since its degree is at most \(2\kappa-2\), it satisfies
\(f(X)=\Psi(X)h(X)\) with \(\deg h\leq\kappa-2\).
At an additional node, coefficient aggregation commutes with
evaluation at \(\ext z\):
\[
\begin{aligned}
    F_{\tau+1,j}(\ext z)
      &=\sum_{\nu=0}^{2^{\tau+1}}
        \langle\ext{\bm l}_{\tau+1},
                 \bm p_{\tau+1,j,\nu}\rangle\ext z^\nu\\
      &=\left\langle\ext{\bm l}_{\tau+1},
             \bm g_{\tau,U}(\ext z,\gamma'_j)
               \circ\bm g_{\tau,D}(\ext z,\gamma'_j)
        \right\rangle.
\end{aligned}
\]
Algorithm~\ref{alg:permutation} therefore recovers \(h\) from
\(\kappa-1\) samples. Reencoding \(\bm C_\tau(Z)\) at each
level keeps the \(X\)-degree of \(\bm g_\tau(Z,X)\) at most
\(\kappa-1\).

Each evaluated relation is proved by weighted sumcheck.
For length-\(n\) child vectors and a tensor weight
\(\ext{\bm l}\), the sumcheck polynomial is:
\begin{equation}
    \widetilde{\ext{\bm l}}(\bm Y)\,
    \widetilde{\ext{\bm u}}(\bm Y)\,
    \widetilde{\ext{\bm v}}(\bm Y).
\label{eq:permutation-hadamard-as-ip}
\end{equation}
It has individual degree at most three.
The two child claims are evaluated at the same point
\(\ext{\bm y}\). Their affine combination with a fresh
\(\ext y_0\) is the MLE claim on the concatenated parent
vector at \(\ext y_0\mathbin\|\ext{\bm y}\).
This preserves the recursive matrix claim
\(\ext v_\tau=\ext{\bm l}_\tau^\top
\bm C_\tau(\ext z)\ext{\bm r}_\tau\).
At level zero, the row MLE weights and the column Lagrange
weights both sum to one. Removing the shift leaves:
\[
    \ext v_0-\ext z
       =\ext{\bm l}_0^\top\bm A\ext{\bm r}_0.
\]
This is the packed-opening interface of
Eq.~\eqref{eq:pcs-opening-relation}. The same argument applies
to \(\bm B\).

\smallskip\noindent\textbf{Soundness.}
For a weighted sumcheck on length-\(n\) vectors,
Eq.~\eqref{eq:permutation-hadamard-as-ip} contributes at most
\(3\log_2n/|\ExtensionField|\).
If at least one of the two terminal child claims is false,
their errors form a nonzero affine polynomial in the fresh
merge challenge \(\ext y_0\). The merged claim is then correct
with probability at most \(1/|\ExtensionField|\).
A transition thus preserves a false claim except with
probability at most \((3\log_2n+1)/|\ExtensionField|\).
The extension-field tree has child lengths
\(1,2,\ldots,m/2\). Summing its errors gives:
\[
    \sum_{i=0}^{\log_2m-1}
       \frac{3i+1}{|\ExtensionField|}
       =\frac{\frac32\log_2m(\log_2m-1)+\log_2m}
                    {|\ExtensionField|}.
\]

At coefficient level \(\tau+1\), a false incoming matrix
claim either fails the scalar check against
\(\ext{\bm u}_{\tau+1}\), or makes at least one column claim
false. In the latter case \(f(X)\) does not vanish on all of
\(\mathcal H\), and \(f(X)-\Psi(X)h(X)\) is nonzero for every
claimed \(h\) of degree at most \(\kappa-2\).
Since the message is fixed before \(\ext x\) is sampled,
this polynomial vanishes at \(\ext x\) with probability at
most \((2\kappa-2)/(|\ExtensionField|-\kappa)\).
The subsequent weighted sumcheck and merge have child length
\(m/2^{\tau+1}\) and error at most
\((3\log_2(m/2^{\tau+1})+1)/|\ExtensionField|\).
Summing over \(\tau=0,\ldots,\log_2\kappa-1\) and adding the
extension-field-tree error gives the stated
\(\epsilon_{\mathsf{GP}}\). This bound is conditional on
authentication of the terminal packed claim.

If the committed vectors are not permutations of one another,
their fingerprint difference is the nonzero polynomial:
\[
    \prod_{i=1}^N(Z+a_i)-\prod_{i=1}^N(Z+b_i).
\]
Its degree is at most \(N-1\), so the true fingerprints
coincide with probability at most
\(\min\{1,(N-1)/|\ExtensionField|\}\).
Otherwise, at least one claimed grand product is false.
A union bound over the two reductions and their terminal
authentication gives:
\[
    \epsilon_{\mathsf{Perm}}
       \leq\min\left\{1,\frac{N-1}{|\ExtensionField|}\right\}
          +2\epsilon_{\mathsf{GP}}+\epsilon_{\mathsf{PCS}}.
\]

\smallskip\noindent\textbf{Efficiency over \(\Field\).}
At level \(\tau\in[\log_2\kappa]\), the coefficient matrix
of \(\bm C_{\tau-1}(Z)\) has \(O(m)\) row-degree pairs.
Evaluating its \(\kappa\)-column encodings at \(\kappa-1\)
additional nodes costs \(O(N\kappa)\,\bb\).
Multiplying the resulting degree-\(2^{\tau-1}\) polynomials
costs \(O(m\kappa2^\tau)\,\bb\).
These products, including construction of the coefficient
tree, sum to \(O(N\kappa)\,\bb\).

The off-domain products and column claims contain \(O(N)\)
coefficients per level, whose aggregation costs \(O(N)\,\be\).
There are \(O(\kappa)\) aggregated polynomials of degree
at most \(2^\tau\), so their subsequent evaluation costs:
\[
    \sum_{\tau=1}^{\log_2\kappa}
        O(\kappa2^\tau)\,\ee
       =O(\kappa^2)\,\ee.
\]
The shared powers of \(\ext z\) require
\(O(\kappa)\,\ee\) once.
Forming the vectors for the short weighted sumchecks costs
\(O(N)\,\be+O(m)\,\ee\) per coefficient level.
Their total sumcheck length, together with the
extension-field tree, is \(O(m)\).
Interpolation and public queries contribute
\(O(\kappa^3)\,\bb\) per level.
Using \(m\geq\kappa\), the total is:
\[
    O(N\kappa\log_2\kappa+\kappa^4)\,\bb.
\]
The \(\kappa^4\) term accounts for evaluation of the
aggregated extension-field polynomials.

\smallskip\noindent\textbf{Intermediate-field implementation.}
For \(q<2\kappa-1\), use the addition and multiplication tables
of Section~\ref{sec:hadamard} over \(\mathbb K\), with the same
unit-cost lookup convention.
Precomputed Lagrange evaluations allow the
\(O(N\kappa)\) encoding operations at each level to be
performed over \(\mathbb K\).
The row-polynomial products likewise use
\(O(N\kappa)\) table operations in total.
The coefficients of \(\bm C_\tau(Z)\) remain in \(\Field\);
only the off-domain encodings and products use \(\mathbb K\).

In a compatible tower basis, a mixed
\(\mathbb K\times\ExtensionField\) product and its
accumulation cost \(O(\kappa)\) in base-field coordinates.
Schoolbook multiplication over \(\mathbb K\), using the
tables for coordinate products, costs
\(O(\kappa^2/t)\) per extension-field multiplication.
The preceding count becomes
\(O\bigl(N\kappa\log_2\kappa+\kappa^4/t\bigr)\,\bb\),
which implies the coarser bound in the proposition.

\smallskip\noindent\textbf{Space.}
One coefficient-tree level contains \(O(N)\) base-field
coefficients. The levels can be reconstructed from
\(\bm C_0(Z)\) during the reverse traversal:
reconstructing level \(\tau\) costs \(O(N2^\tau)\,\bb\),
and these costs sum to \(O(N\kappa)\,\bb\).
Thus the prover need not retain all coefficient levels.
Processing one additional node at a time uses
\(O(mt)\subseteq O(N)\) coordinates for its row polynomials.
The aggregated polynomials use \(O(\kappa^2)\)
extension-field coefficients, or \(O(\kappa^3)\)
base-field coordinates. The extension-field tree and
sumcheck tables use \(O(m\kappa)=O(N)\) coordinates.
Including interpolation data and the intermediate-field
tables, peak memory is \(O(N+\kappa^3+tq^{2t})\)
base-field coordinates.
When \(q\geq2\kappa-1\), no field tables are needed, giving
\(O(N+\kappa^3)\) base-field coordinates.
The time bounds exclude preprocessing, commitments, and
terminal openings.
\end{proof}

\subsection{Proof of Proposition~\ref{proposition:packed-pcs}}
\label{appx:proof-pcs}

\begin{proof}
\noindent\textbf{Binding and completeness.}
Each entry of \(\bm A\ext{\bm\beta}\) uniquely determines
the corresponding row of \(\bm A\). This injective encoding
preserves binding and, when provided by the underlying PCS,
commitment hiding.

For an honest prover, \(\bm U=\bm\Xi^\top\bm A\), so
Eq.~\eqref{eq:pcs-encoded-matrix} and its random projection hold.
We show that the verifier also computes the public terminal value
exactly.

By Eq.~\eqref{eq:pcs-multiplication-matrix}, multiplication by
\(\ext z\) acts on a column of base-field coordinates through
\(\bm M_{\ext z}\). This map is \(\Field\)-linear in \(\ext z\),
and \(\bm M_{1}=\bm I_\kappa\).
We prove by backward induction that
\(\bm W_j=\bm Y_j^\top\bm\Xi_j\).
Since \(\ext\beta_1=1\), the empty suffix has coordinate row
\(\bm e_1^\top\), giving
\(\bm Y_{\ell+1}^\top\bm\Xi_{\ell+1}
=\bm e_1\bm e_1^\top=\bm W_{\ell+1}\).
For the induction step, apply the block identities in
Eq.~\eqref{eq:pcs-query-blocks} and the induction hypothesis:
\[
\begin{aligned}
    \bm Y_j^\top\bm\Xi_j
    &=\bm M_{1-\ext y_j}\bm W_{j+1}
        \bm M_{\ext\xi^{(j)}_0}^\top\\
    &\quad{}+\bm M_{\ext y_j}\bm W_{j+1}
        \bm M_{\ext\xi^{(j)}_1}^\top\\
    &=\bm W_{j+1}\bm M_{\ext\xi^{(j)}_0}^\top\\
    &\quad{}+\bm M_{\ext y_j}\bm W_{j+1}
        \bm M_{\ext\xi^{(j)}_1-\ext\xi^{(j)}_0}^\top\\
    &=\bm W_j.
\end{aligned}
\]
The second equality uses linearity of \(\bm M\); the last
is Eq.~\eqref{eq:pcs-query-recurrence}.
Taking \(j=1\) gives \(\bm W=\bm W_1=\bm Y^\top\bm\Xi\)
and hence:
\[
    \ext{\bm\beta}^\top\bm W\ext{\bm r}
    =(\bm Y\ext{\bm\beta})^\top\bm\Xi\ext{\bm r}
    =\tensorspan{\ext{\bm y}}^\top\bm\Xi\ext{\bm r}.
\]
Thus Eq.~\eqref{eq:pcs-query-matrix} checks the query-side
terminal value exactly. Completeness follows from
Fact~\ref{fact:ip} and the underlying PCS opening.
For \(m=1\), the initialization suffices, with no recurrence
or sumcheck rounds.

\smallskip\noindent\textbf{Soundness.}
Fix the matrix \(\bm A\) bound by the commitment and suppose that
Eq.~\eqref{eq:pcs-opening-relation} is false.
If the first scalar check accepts, then
\(\ext{\bm u}^{\top}\neq\ext{\bm\xi}^{\top}\bm A\), hence
\(\bm U\neq\bm\Xi^\top\bm A\).
Injectivity of basis packing makes
\((\bm U-\bm\Xi^\top\bm A)\ext{\bm\beta}\) a nonzero
column of \(\ExtensionField^\kappa\), fixed before
\(\ext{\bm r}\) is sampled. Its inner product with the uniform
\(\ext{\bm r}\) vanishes with probability
\(1/|\ExtensionField|\). Outside this event,
Eq.~\eqref{eq:pcs-random-opening} is false.

Except with probability \(2\ell/|\ExtensionField|\),
Fact~\ref{fact:ip} then leaves a false terminal claim in
Eq.~\eqref{eq:pcs-terminal-claims}.
The recurrence computes \(\bm W\) from the public tensor factors
and the terminal point, so a false query-side claim is rejected
deterministically. If that check passes, the witness-side claim
must be false. The underlying PCS accepts this adaptively chosen
false MLE claim with probability at most
\(\epsilon_{\mathsf{PCS}}\).
A union bound gives Eq.~\eqref{eq:pcs-soundness}.

\smallskip\noindent\textbf{Prover complexity.}
Expanding \(\ext{\bm\xi}\) and reading its basis coordinates
gives \(\bm\Xi\) in \(O(m)\) extension-field operations.
The matrix product \(\bm U=\bm\Xi^\top\bm A\) takes
\(\kappa N\) base-field multiplications. Its columns give the
coordinates of \(\ext{\bm u}^{\top}=\ext{\bm\beta}^{\top}\bm U\),
so packing requires only a change of representation.

Computing \(\bm\Xi\ext{\bm r}\) takes \(N\) mixed-field
multiplications; the extension-field inner-product reduction takes
\(4m+O(\ell)\) extension-field multiplications.
Packing the rows of \(\bm U\) and taking their inner product
with \(\ext{\bm r}\) computes \(\ext c\) using
\(\kappa\) extension-field multiplications.
The prover need not construct \(\bm Y\) or \(\bm W\).
These costs give Eq.~\eqref{eq:pcs-prover-cost}.
Under schoolbook arithmetic they are
\(O(\kappa N+\kappa^3)\,\bb\), which is
\(O(\kappa N)\,\bb\) when \(m\geq\kappa\).

The matrices \(\bm A,\bm\Xi\) occupy \(O(N)\) base-field
coordinates, and \(\bm U\) and the length-\(\kappa\)
extension-field vectors occupy \(O(\kappa^2)\).
A constant number of length-\(m\) extension-field buffers also
occupy \(O(m\kappa)=O(N)\) base-field coordinates.
Thus peak prover memory, excluding the underlying PCS workspace,
is \(O(N+\kappa^2)\), or \(O(N)\) when \(m\geq\kappa\).

\smallskip\noindent\textbf{Verifier complexity and communication.}
The verifier implements Eq.~\eqref{eq:pcs-query-recurrence}
without forming the multiplication matrices.
Right multiplication by \(\bm M_{\ext z}^\top\) scales each
row, interpreted in the fixed basis as an extension-field element,
by \(\ext z\). Left multiplication by \(\bm M_{\ext z}\)
scales each column in the same way. Each scaling therefore costs
\(\kappa\) extension-field multiplications.
The two terms in each recurrence step require three such scalings,
followed by \(\kappa^2\) base-field additions.
Over \(\ell\) steps, the cost is at most
\(3\kappa\ell\) extension-field multiplications and
\(O(\kappa^2\ell)\) base-field operations, including coordinate
rearrangement. A constant number of intermediate matrices use
\(O(\kappa^2)\) additional base-field coordinates.

The initial scalar check costs \(\kappa\) extension-field
multiplications. To compute \(\ext c\) and the expected terminal
value \(\ext w\), the verifier packs the rows of \(\bm U\) and
the columns of \(\bm W\), respectively, and takes one
length-\(\kappa\) extension-field inner product per target.
Packing is coordinate reinterpretation, so these targets
cost \(O(\kappa)\) extension-field multiplications and
\(O(\kappa^2)\) coordinate operations in total.
The inner-product verifier costs \(O(1+\ell)\) extension-field
operations, including its terminal scalar check.
This gives Eq.~\eqref{eq:pcs-verifier-cost}.

The prover sends \(\kappa\) extension-field elements for
\(\ext{\bm u}\), followed by the \(2\ell+2\) extension-field
elements of Fact~\ref{fact:ip}, including both terminal values.
The verifier computes \(\bm W\) locally, so no matrix message or
base-field proof is sent. Excluding the underlying opening, total
prover-to-verifier communication is therefore
\(\kappa+2\ell+2\) extension-field elements.
\end{proof}

\subsection{Proof of Lemma~\ref{eq:permutation}}
\label{appx:proof-copy-constraint}

\begin{proof}
The radix-\(\underline q\) representations of \(0,\ldots,L-1\)
use at most \(\lceil\log_{\underline q}L\rceil\) digits.
The digit injection preserves their distinctness.
Since \(\theta_1,\ldots,\theta_d\) are linearly independent
over \(\Field\), and the last coordinate is reserved for the
witness value:
\[
    a_i=b_j
    \quad\Longleftrightarrow\quad
    \bigl(\sigma(i)=j\ \text{and}\ w_i=w_j\bigr).
\]
Every label occurs exactly once in each operand.
Thus equality of the multisets forces
\(a_i=b_{\sigma(i)}\), and hence \(w_i=w_{\sigma(i)}\),
for all \(i\in[L]\).
Conversely, these copy equalities give
\(a_i=b_{\sigma(i)}\), proving the permutation relation.
The encoding is deterministic and injective, so it adds no
soundness error. Appending the same zeros to both operands
preserves the equivalence by multiset cancellation.
\end{proof}

\subsection{Proof of Proposition~\ref{proposition:position-encoding-evaluation}}
\label{appx:proof-position-encoding}

\begin{proof}
For the chosen digit injection, each label is the sum of the
contributions of its nonzero position bits.
The recursion and interval decomposition in
Appendix~\ref{appx:mleiota} therefore give the ordinary MLE
evaluation in \(O(\log L)\) operations.

For a packed query, split the position bits into column bits
and row bits in column-major order.
The unpadded label is the sum of a column contribution
\(\bm\iota_{\mathsf{col}}\) and a row contribution
\(\bm\iota_{\mathsf{row}}\).
For arbitrary column weights \(\ext{\bm r}\) and tensor row
weights \(\ext{\bm l}\), the full-table contraction is:
\[
    \left\langle\bm\iota_{\mathsf{col}},\ext{\bm r}\right\rangle
       \left(\bm1^\top\ext{\bm l}\right)
    +
    \left(\bm1^\top\ext{\bm r}\right)
       \left\langle\bm\iota_{\mathsf{row}},\ext{\bm l}\right\rangle.
\]
Recursive doubling generates all \(\kappa_{\mathbb K}\)
column contributions in \(O(\kappa_{\mathbb K})\) additions.
For each row-bit suffix, compute both its total tensor weight
and its label-weighted sum.
Splitting off one two-entry tensor factor updates both
quantities in constant time, including when its entries do
not sum to one.

Zero padding keeps the first \(L\) positions.
In column-major order, these consist of complete columns
and at most one partially filled column.
All complete columns use the same two full-row contractions.
For the partial column, the binary interval decomposition
computes the two corresponding prefix contractions using the
precomputed suffix quantities.
One scan of the prefix length costs \(O(\log L)\) operations.
Combining the column contributions then costs
\(O(\kappa_{\mathbb K})\) operations.
The total is \(O(\kappa_{\mathbb K}+\log L)\), with the same
bound on working field elements and no division.
\end{proof}

\subsection{Proof of Theorem~\ref{theorem:single-instance-packing}}
\label{appx:proof-single-instance}

\begin{proof}
\noindent\textbf{Completeness.}
For a satisfying assignment, the five scalar claims in
Eq.~\eqref{eq:plonkWeightedHadamard1} and
Eq.~\eqref{eq:plonkWeightedHadamard2} sum to zero.
Four packed Hadamard proofs reduce the selector--wire product
claims to evaluations of the selectors and
\(\bm w_L,\bm w_R,\bm w_M,\bm w_O\).
The fifth authenticates the remaining evaluation of
\(\bm w_M=\bm w_L\circ\bm w_R\).
All terminal queries have arbitrary column weights and tensor
row weights, as required by Section~\ref{sec:pcs}.

Lemma~\ref{eq:permutation} proves the copy equivalence.
The two operands have length \(L=3N+n_{\mathsf{io}}\leq4N\);
append zeros to reach length \(4N\).
Their grand products then share the padding factor
\(\ext z^{4N-L}\), which is nonzero because
\(\ext z\sample\ExtensionField\setminus\{0\}\).
The two grand-product reductions therefore prove exactly the
original permutation relation.
Their packing width \(\kappa_{\mathbb K}=\kappa/d\) and
row count \(4N/\kappa_{\mathbb K}=4md\) are powers of two,
with \(4md\geq\kappa_{\mathbb K}\geq2\).
Thus the parameters satisfy the interface of
Section~\ref{sec:permutation}, and every honest check accepts.

\smallskip\noindent\textbf{Terminal copy queries.}
Each copy terminal is a weighted evaluation of a padded operand,
with \(\kappa_{\mathbb K}\) arbitrary column weights and
\(4N/\kappa_{\mathbb K}\) tensor row weights.
By linearity, the two terminals involve
\(\bm\varsigma,\bm\iota\), and \(\bm w\).
Restricting to one length-\(N\) wire segment fixes the first
two position bits.
The last \(\log_2m\) bits remain tensor row factors.
Absorb all other free factors into the column weights to obtain
a length-\(\kappa\) column query on the original
\(m\times\kappa\) wire or digit matrix.
This changes the query representation, not the committed data.

Proposition~\ref{proposition:position-encoding-evaluation}
computes the identity-label contribution in
\(O(\kappa_{\mathbb K}+\log L)\) extension-field operations.
The verifier evaluates the public-I/O values and their permuted
labels directly. Traversing their contiguous interval and
reusing tensor prefixes costs
\(O(\kappa_{\mathbb K}+n_{\mathsf{io}}+\log L)\) operations.
No commitment to either contribution is needed.

\smallskip\noindent\textbf{Batched terminal openings.}
Write \(h=\lceil\log_{\underline q}L\rceil\).
The three fixed-basis wire packings form one witness vector
of logical length \(3m\), committed before any challenge.
The public vector concatenates the five selector packings
and the three wire-segment restrictions of each of the \(h\)
permuted-position digit tables.
Its logical length is \(m_{\mathsf{pub}}=(5+3h)m\).
The public commitment, circuit tables, and packing polynomials
are preprocessed.

The gate argument leaves five selector and five wire queries.
The two copy terminals add six wire queries and \(3h\)
permuted-digit queries, giving \(16+3h\) packed queries.
Apply the column-projection step of Section~\ref{sec:pcs}
to each query and check its claimed scalar value.
Fix all column messages before sampling independent projection
vectors for the different queries.
Within each group, zero-extend the projected queries to their
committed blocks and sum the resulting inner-product claims.
One invocation of Fact~\ref{fact:ip} per group reduces their
sum to one MLE opening of the concatenated vector.
The verifier evaluates the combined public query exactly,
using the recurrence of Section~\ref{sec:pcs} on each block
and retaining its block-prefix weight.
Thus there are two underlying PCS openings.

If a column message is false, its coordinate error satisfies
\((\bm U-\bm\Xi^\top\bm A)\ext{\bm\beta}\neq\bm0\).
Consequently, a group containing a false claim has a nonzero
linear error polynomial in its independent projection vectors.
It vanishes with probability at most
\(1/|\ExtensionField|\) per group.
The two inner-product reductions add
\(O(\log N)/|\ExtensionField|\) error.
Here the padded group lengths have logarithm \(O(\log N)\),
since \(h<d\leq\kappa/2\).

\smallskip\noindent\textbf{Soundness.}
Condition on binding of the witness commitment.
Its fixed-basis packing determines all three wire vectors
and the virtual product before any challenge.
The selectors and position tables are fixed at setup.

A nonzero gate residual vanishes at the random MLE point with
probability at most \(\log_2N/|\ExtensionField|\).
Otherwise an accepting scalar sum contains a false claim.
Acceptance then requires an error in one of the five Hadamard
reductions or in terminal authentication.
For the nested product claim, its value is fixed before the
fresh challenges of the fifth reduction.

If a copy constraint fails, the unpadded fingerprint difference
is nonzero and has degree at most \(L-1\).
At the nonzero challenge \(\ext z\), its padded version vanishes
if and only if the unpadded difference does.
The fingerprint error is therefore at most
\((L-1)/(|\ExtensionField|-1)\).
Let \(\epsilon_{\mathsf{GP}}\) be the bound of
Proposition~\ref{proposition:permutation-complexity},
instantiated with length \(4N\), packing width
\(\kappa_{\mathbb K}\), and base field \(\mathbb K\).
A union bound gives:
\[
\begin{aligned}
    \epsilon_{\mathsf{single}}\leq{}&
    \epsilon_{\mathsf{PCS}}
    +\frac{L-1}{|\ExtensionField|-1}
    +2\epsilon_{\mathsf{GP}}\\
    &+5\left(
        \frac{2\kappa-2}{|\ExtensionField|-\kappa}
        +\frac{2\log_2N}{|\ExtensionField|}
      \right)
    +O\left(\frac{\log N}{|\ExtensionField|}\right).
\end{aligned}
\]
The last term includes the initial gate projection and the
two batched opening reductions.
The grand-product contribution is
\(O((\kappa_{\mathbb K}\log_2\kappa_{\mathbb K}
+(\log_2N)^2)/|\ExtensionField|)\).
Since \(N=m\kappa\geq\kappa^2\), \(L\leq4N\), and
\((\log_2N)^2=O(N)\), this proves
Eq.~\eqref{eq:together-soundness}.
The term \(\epsilon_{\mathsf{PCS}}\) includes joint binding
and the two adaptive underlying openings.

\smallskip\noindent\textbf{Prover time.}
The five Hadamard reductions cost \(O(N\kappa)\bb\).
For \(q\geq2\kappa-1\), interpolation uses base-field nodes.
If \(\kappa<q<2\kappa-1\), the intermediate-field variant
uses degree-two interpolation-field tables; its online cost
is unchanged and its table storage is \(O(q^4)=O(\kappa^4)\).
The power-of-two \(\kappa\) permits this degree-two subfield.
The \(16+3h\) column projections and the two grouped
inner-product reductions cost \(O(N\kappa h)\bb\).

Choose polynomial bases along \(\Field\subset\mathbb K
\subset\ExtensionField\) for the copy computation.
A \(\mathbb K\)-multiplication costs \(O(d^2)\bb\).
The two grand-product reductions therefore cost:
\[
    O\left(
       N\kappa_{\mathbb K}\log_2\kappa_{\mathbb K}
       +\kappa_{\mathbb K}^4
    \right)d^2\bb
    =
    O\left(
       N\kappa d\log_2\kappa_{\mathbb K}
       +\frac{\kappa^4}{d^2}
    \right)\bb.
\]
The \(\log_2\kappa_{\mathbb K}\) factor counts the packed
coefficient levels; evaluating their aggregated polynomials
accounts for the additive term.
The least-power-of-two choice gives
\(h+1\leq d<2(h+1)\).
Since \(\log\underline q=\Theta(\log q)\) and \(L=\Theta(N)\),
both \(h\) and \(d\) are \(\Theta(\log_qN)\).
Combining these contributions gives the stated
\(O(N\kappa\log_qN\log_2\kappa)\bb+O(\kappa^4)\) bound.
The copy computation stays in its tower basis; only interface
queries and messages are converted to the fixed commitment basis.
These conversions cost \(O(\kappa^2)\bb\) per extension-field
element and are covered by the same bounds.

\smallskip\noindent\textbf{Verifier time and communication.}
The \(O(h)\) packed column checks and combined query evaluations
cost \(O(\kappa h\log_2N)\ee\).
The two grand-product reductions add
\(O(\kappa_{\mathbb K}\log_2\kappa_{\mathbb K}
+(\log_2N)^2)\ee\).
Direct identity-label and public-I/O evaluations contribute
\(O(\kappa_{\mathbb K}+n_{\mathsf{io}}+\log N)\ee\).
These costs give the verifier bound.

The packed column messages contain \(O(\kappa h)\)
extension-field elements.
The two grand-product transcripts contain
\(O(\kappa_{\mathbb K}\log_2\kappa_{\mathbb K}
+(\log_2N)^2)\) such elements.
The remaining Hadamard and batched inner-product transcripts
are absorbed by these terms.
Using \(h=O(\log_qN)\) and
\(\kappa_{\mathbb K}\leq\kappa\) gives the stated proof size.

\smallskip\noindent\textbf{Prover memory and PCS costs.}
The grand-product workspace uses
\(O(dN+d\kappa_{\mathbb K}^3)
=O(dN+\kappa^3/d^2)\) base-field coordinates.
The witnesses, circuit tables, and grouped opening buffers
occupy \(O(Nh)\) coordinates.
Sequential execution reuses the gate and copy workspaces.
Including the optional interpolation tables and public I/O
gives \(O(N\log_qN)\Field+O(\kappa^4+n_{\mathsf{io}})\).

The four PCS functions bound aggregate resources for one online
witness commitment and one opening of each of the two vectors,
of logical lengths \(3m\) and \(m_{\mathsf{pub}}\).
They include domain padding and both retained commitment states;
public commitment generation is preprocessed.
Their total argument is therefore \(3m+m_{\mathsf{pub}}\).
This proves Eq.~\eqref{eq:together-single-efficiency}.
Unlabeled time and memory remainders count base-field operations
and base-field elements, respectively.
\end{proof}

\subsection{Proof of Theorem~\ref{theorem:multiple-instance-packing}}
\label{appx:proof-multiple-instance}

\begin{proof}
\noindent\textbf{Completeness.}
Let \(\bm R\in\Field^{N\times\kappa}\) contain the gate
residuals of the \(\kappa\) instances as columns.
The left-hand side of Eq.~\eqref{eq:together-multiple-gates}
is \(\bm R\ext{\bm\beta}\), including the term
\(\bm q_C\bm1_\kappa^\top\ext{\bm\beta}\).
Linear independence of the basis gives
\(\bm R\ext{\bm\beta}=\bm0_N\) if and only if
\(\bm R=\bm0_{N\times\kappa}\).
Similarly, packing commutes with the common row permutation,
so \(\sigma(\ext{\bm w})=\ext{\bm w}\) is equivalent to
the copy constraints in every instance.
This equivalence includes the packed public I/O.
Eq.~\eqref{eq:together-multiple-product-opening} authenticates
the virtual product vector from
Eq.~\eqref{eq:together-multiple-wires}.
Thus all scalar checks and tree relations accept for an honest prover.

\smallskip\noindent\textbf{Public packing.}
Let \(h=\lceil\log_{\underline q}L\rceil\).
Form \(\bm Q\in\Field^{N\times\kappa}\) with first five columns
\(\bm q_L,\bm q_R,\bm q_M,\bm q_O,\bm q_C\).
The next three blocks of \(h\) columns contain the position digits
\(\bm\sigma_{(1)},\ldots,\bm\sigma_{(h)}\), restricted to the
three length-\(N\) wire segments, respectively.
Fill the remaining columns with zeros.
The capacity condition \(5+3h\leq\kappa\) therefore gives:
\[
    \bm Q\ext{\bm\beta}\in\ExtensionField^N.
\]
A selector query uses its coordinate unit vector as the column
weight in Eq.~\eqref{eq:pcs-opening-relation}.
For the permuted-label query on segment \(\ell\in[3]\), use
\(\ext\rho_{5+(\ell-1)h+j}=\ext\beta_j\) for \(j\in[h]\),
and zero elsewhere.
Thus five selector queries and three permuted-label queries
are weighted evaluations of the same public matrix.
The row weights are tensor vectors, so Section~\ref{sec:pcs}
applies to all eight queries.
The public commitment is preprocessed.
The identity-label MLE is evaluated directly by
Proposition~\ref{proposition:position-encoding-evaluation};
it is not committed.

\smallskip\noindent\textbf{Gate soundness.}
Condition on binding of the witness commitment.
Its three length-\(N\) blocks have unique matrices of
\(\Field\)-coordinates in the fixed basis.
Thus the gate residual and \(\ext{\bm w}_M\) are determined
by the committed witnesses.
If some gate is invalid, the packed residual is nonzero.
Its random length-\(N\) tensor projection vanishes with
probability at most \(\log_2N/|\ExtensionField|\).

For a fixed projection point \(\ext{\bm x}\), the gate
sumcheck polynomial is the product of the multilinear equality
weight with the residual expression formed from selector and wire
MLEs. Each selector--wire product has individual degree two,
and multiplication by the equality weight raises this to three.
The constant-selector term has lower degree, giving sumcheck
error at most \(3\log_2N/|\ExtensionField|\).
The remaining virtual-product evaluation is checked by one
packed Hadamard proof on matrices of size \(N\times\kappa\).
Since the total number of entries is \(N\kappa\), its
reduction error is bounded by:
\[
    \frac{2\kappa-2}{|\ExtensionField|-\kappa}
       +\frac{2\log_2(N\kappa)}{|\ExtensionField|}.
\]
All terminal queries are authenticated by the batched openings below.
The bound applies jointly to all \(\kappa\) instances because
packing is injective.

\smallskip\noindent\textbf{Copy fingerprint.}
Suppose that the packed assignment violates a copy constraint.
Consider the bivariate difference polynomial:
\[
    P(Y,Z)=
       \prod_{i=1}^{L}(Z+\iota_i+Y\ext w_i)
       -\prod_{i=1}^{L}(Z+\varsigma_i+Y\ext w_i).
\]
If \(P\) were zero, unique factorization in
\(\ExtensionField[Y,Z]\) would equate the two multisets of
linear factors, each monic in \(Z\).
Since the labels are distinct and
\(\varsigma_i=\iota_{\sigma(i)}\), matching the constant
coefficients pairs position \(i\) with \(\sigma^{-1}(i)\).
Matching the \(Y\)-coefficients then forces
\(\ext w_i=\ext w_{\sigma^{-1}(i)}\), contradicting the copy
violation. Hence \(P\ne0\).

Both products have homogeneous degree-\(L\) part
\(\prod_{i=1}^{L}(Z+Y\ext w_i)\), so \(P\) has total degree
at most \(L-1\).
The independent uniform challenges \(\ext y,\ext z\)
therefore satisfy \(P(\ext y,\ext z)=0\) with probability at
most \(\min\{1,(L-1)/|\ExtensionField|\}\).
Padding each product-tree input with ones does not change
this fingerprint.

\smallskip\noindent\textbf{Product-tree soundness.}
Set \(\mu=\lceil\log_2L\rceil=\log_2(4N)\).
For one grand product, unfold the binary tree from its root
to the padded length-\(4N\) input vector.
At a level with \(2^\ell\) parent entries, where
\(\ell\in\{0,\ldots,\mu-1\}\), let
\(\ext{\bm x}\in\ExtensionField^\ell\) be the current
evaluation point.
The weighted product identity uses the polynomial:
\[
    \left\langle
       \tensorspan{\ext{\bm x}},\tensorspan{\bm X}
    \right\rangle
    \left\langle\ext{\bm c}_U,\tensorspan{\bm X}\right\rangle
    \left\langle\ext{\bm c}_D,\tensorspan{\bm X}\right\rangle,
\]
where \(\ext{\bm c}_U,\ext{\bm c}_D\) are the two child subvectors
and \(\bm X\) denotes \(\ell\) formal variables.
Its degree in each variable is at most three.
Sumcheck reduces its Boolean sum to child evaluations
\(\ext v_U,\ext v_D\) at one point \(\ext{\bm y}\), with
error at most \(3\ell/|\ExtensionField|\).
After both values are fixed, sample a fresh
\(\ext y_0\sample\ExtensionField\) and form the next claim:
\[
    (1-\ext y_0)\ext v_U+\ext y_0\ext v_D
       =
    \left\langle
       \ext{\bm c}_U\mathbin{\|}\ext{\bm c}_D,
       \tensorspan{(\ext y_0,\ext{\bm y})}
    \right\rangle.
\]
If either child claim is false, this merge accepts it with
probability at most \(1/|\ExtensionField|\).
Thus the two trees have total reduction error at most:
\[
    2\sum_{\ell=0}^{\mu-1}\frac{3\ell+1}{|\ExtensionField|}
       =\frac{3\mu(\mu-1)+2\mu}{|\ExtensionField|}.
\]
The \(\ell=0\) level has no sumcheck rounds but still uses
a fresh merge challenge.

Each tree ends in a single MLE claim on its padded input.
At its terminal point, the known shift and padding contribute:
\[
    \ext z\sum_{i=1}^{L}
       \bigl(\tensorspan{\ext{\bm x}}\bigr)_i
    +\sum_{i=L+1}^{4N}
       \bigl(\tensorspan{\ext{\bm x}}\bigr)_i
    =
    1+(\ext z-1)\sum_{i=1}^{L}
       \bigl(\tensorspan{\ext{\bm x}}\bigr)_i.
\]
The prefix sum takes \(O(\log N)\) operations by a binary
interval decomposition.
Each length-\(N\) wire segment fixes two prefix bits.
Retain their tensor weight and query the suffix MLE.
The public-I/O values and permuted labels at the remaining
\(n_{\mathsf{io}}\) positions are evaluated directly.
The identity-label MLE is computed by Appendix~\ref{appx:mleiota}.
Consequently, the two copy terminals leave six ordinary wire
queries and three weighted queries to the public matrix.

\smallskip\noindent\textbf{Batched openings and composition.}
There are nine ordinary wire queries: three from the gate check
and six from the copy check.
The two Hadamard terminals and eight public queries give
ten packed queries.
Apply the column-projection step of Section~\ref{sec:pcs}
to each packed query.
Fix all column messages and ordinary claimed values before
sampling independent projection vectors for the packed queries
and independent scalar weights for the ordinary queries.
Extend each query by zeros outside its committed block, and
sum the resulting inner-product claims within each group.

If an ordinary claim is false, its scalar weight has a nonzero
error coefficient.
If a packed column message is false and its scalar check passes,
injectivity gives
\((\bm U-\bm\Xi^\top\bm A)\ext{\bm\beta}\ne\bm0\).
Thus a group containing a false claim has a nonzero linear
error polynomial in its independent batching challenges.
Cancellation occurs with probability at most
\(1/|\ExtensionField|\) per group.

The witness concatenation has logical length \(3N\), padded
to \(4N\); the public packing has length \(N\).
One invocation of Fact~\ref{fact:ip} per group therefore adds
\(2\log_2(4N)/|\ExtensionField|\) and
\(2\log_2N/|\ExtensionField|\) error, respectively.
The verifier computes the combined query-side evaluation in
each group exactly: use the recurrence of Section~\ref{sec:pcs}
for packed queries and tensor contractions for ordinary
queries, retaining the witness block-prefix weights.
This leaves two underlying PCS openings.
Their joint binding and opening error is
\(\epsilon_{\mathsf{PCS}}\).

Combining these bounds gives:
\[
\begin{aligned}
    \epsilon_{\mathsf{multi}}
    \leq{}&\epsilon_{\mathsf{PCS}}
       +\frac{2\kappa-2}{|\ExtensionField|-\kappa}\\
       &+\frac{
          L-1+3\mu(\mu-1)+2\mu
          +8\log_2N+2\log_2(N\kappa)+6
       }{|\ExtensionField|}.
\end{aligned}
\]
Here the constant \(6\) includes the two aggregation checks
and the extra prefix rounds for the length-\(4N\) witness.
Since \(L\leq4N\), \(\kappa\leq N\), and
\((\log_2N)^2=O(N)\), this proves
Eq.~\eqref{eq:together-multiple-soundness}.

\smallskip\noindent\textbf{Efficiency.}
Each wire matrix has \(N\kappa\) entries.
Computing \(\bm W_L\circ\bm W_R\) costs \(O(N\kappa)\bb\).
The Hadamard reduction and ten packed column projections each
cost \(O(N\kappa^2+\kappa^3)\bb\), which is
\(O(N\kappa^2)\bb\) because \(N\geq\kappa\).
The assumption \(q\geq2\kappa-1\) supplies the interpolation
nodes over \(\Field\), so no intermediate-field table is needed.
The gate sumcheck and the two final inner-product reductions
use \(O(N)\) extension-field operations.
The product-tree layer lengths sum to \(O(N)\), and their
sumchecks cost \(O(N+\mu^2)\) extension-field operations.
Under schoolbook arithmetic, these contributions give
\(O(N\kappa^2)\bb\) total algebraic prover time.

The packed Hadamard check uses \(O(\kappa)\ee\) verifier time,
and the ten column projections and terminal query contractions
use \(O(\kappa\log_2N)\ee\).
The gate, tree, and final inner-product checks add
\(O((\log_2N)^2)\ee\).
Packing the public I/O reads \(O(\kappa n_{\mathsf{io}})\)
base-field coordinates, and its terminal contractions and
the direct identity-label evaluations take
\(O(n_{\mathsf{io}}+\log N)\) extension-field operations.
This gives the stated verifier time.

The ten column messages contain \(O(\kappa)\)
extension-field elements in total.
The Hadamard proof, gate sumcheck, tree reductions, and final
inner-product proofs add
\(O(\kappa+(\log_2N)^2)\) extension-field elements,
giving the proof-size bound.
The witnesses and public matrix occupy \(O(N\kappa)\)
base-field coordinates.
Sequential execution retains \(O(N)\) extension-field
tree and sumcheck entries and \(O(\kappa^2)\) coordinates
for the column messages.
Thus algebraic prover memory is \(O(N\kappa)\Field\).

The online PCS costs cover one witness commitment and two
openings on vectors of logical lengths \(3N\) and \(N\).
Their aggregate argument is therefore \(4N\), with the required
witness padding included in the PCS cost functions.
Public commitment generation is part of circuit preprocessing;
both retained commitment states count toward PCS memory.
These calculations prove
Eq.~\eqref{eq:together-multiple-efficiency}.
Dividing the algebraic prover time and peak memory by
\(\kappa\) gives the amortized per-instance bounds.
\end{proof}


\end{document}
